% Generated mechanically from frozen input inputs/p12/ja/stacks-geometry.tex.
% Do not edit this generated file; edit the bound locale source instead.

% OK, start here.
%

\setcounter{chapter}{106}
\chapter{代数スタックの幾何学}



\phantomsection
\label{stacks-geometry-section-phantom}


\section{序論}
\label{stacks-geometry-section-introduction}

\noindent
本章では，代数スタックのいくつかの幾何学的性質を論じる．
Sections \ref{stacks-geometry-section-multiplicities} および
\ref{stacks-geometry-section-dimension-of-algebraic-stacks} の初期版は
Matthew Emerton と Toby Gee によって書かれたものであり，
その原形は \cite{Emerton-Gee-dim} に収められている．






\section{半普遍変形環}
\label{stacks-geometry-section-versal}

\noindent
本節では，局所 Noether 的基底上有限型な代数スタックについて，
変形環と局所環との関係を明らかにする．

\begin{situation}
\label{stacks-geometry-situation-versal}
以下，$\mathcal{X}$ は局所 Noether 的スキーム $S$ 上
局所有限型な代数スタックとする．
\end{situation}

\noindent
まず定義を与える．

\begin{definition}
\label{stacks-geometry-definition-versal-ring-at-x}
Situation \ref{stacks-geometry-situation-versal} において，
$x_0 : \Spec(k) \to \mathcal{X}$ を射とし，
$k$ は $S$ 上有限型な体とする．
{\it $x_0$ における $\mathcal{X}$ の半普遍変形環}とは，
剰余体が $k$ である完備 Noether 的局所 $S$-代数 $A$ であって，
Artin's Axioms, Definition
\ref{artin-definition-versal-formal-object}
の意味での半普遍形式対象 $(A, \xi_n, f_n)$ が存在し，
$\xi_1 \cong x_0$（$2$-同型）となるものをいう．
\end{definition}

\noindent
半普遍変形環が存在し，平滑因子を除いて一意であることを示したい．
そのために，Artin's Axioms, Section
\ref{artin-section-predeformation-categories}
の前変形圏を用いる．この状況では，これらは常に変形圏である．

\begin{lemma}
\label{stacks-geometry-lemma-deformation-category}
Situation \ref{stacks-geometry-situation-versal} において，
$x_0 : \Spec(k) \to \mathcal{X}$ を射とし，
$k$ は $S$ 上有限型な体とする．このとき
$\mathcal{F}_{\mathcal{X}, k, x_0}$ は変形圏であり，
$T\mathcal{F}_{\mathcal{X}, k, x_0}$ および
$\text{Inf}(\mathcal{F}_{\mathcal{X}, k, x_0})$
は有限次元 $k$-ベクトル空間である．
\end{lemma}

\begin{proof}
$\Spec(k) \to S$ が経由するアフィン開部分
$\Spec(\Lambda) \subset S$ を選ぶ．
Artin's Axioms, Section \ref{artin-section-predeformation-categories}
により，圏 $\mathcal{C}_\Lambda$ 上の前変形圏
$\mathcal{F}_{\mathcal{X}, k, x_0}$ が得られる．
（前掲箇所で指摘されているように，この圏は
$\Lambda$ の選択には依存せず，射 $\Spec(k) \to S$ のみに依存する．）
Artin's Axioms, Lemmas
\ref{artin-lemma-deformation-category} および
\ref{artin-lemma-algebraic-stack-RS}
により，$\mathcal{F}_{\mathcal{X}, k, x_0}$ は実際に変形圏である．
Artin's Axioms, Lemma \ref{artin-lemma-finite-dimension}
により，$T\mathcal{F}_{\mathcal{X}, k, x_0}$ および
$\text{Inf}(\mathcal{F}_{\mathcal{X}, k, x_0})$
は有限次元 $k$-ベクトル空間である．
\end{proof}

\begin{lemma}
\label{stacks-geometry-lemma-versal-ring}
Situation \ref{stacks-geometry-situation-versal} において，
$x_0 : \Spec(k) \to \mathcal{X}$ を射とし，
$k$ は $S$ 上有限型な体とする．このとき
$x_0$ における $\mathcal{X}$ の半普遍変形環が存在する．
そのような環の対 $A$, $A'$ が与えられれば，ある $r$ に対して
$S$-代数として
$A \cong A'[[t_1, \ldots, t_r]]$
または $A' \cong A[[t_1, \ldots, t_r]]$ である．
\end{lemma}

\begin{proof}
存在は Lemma \ref{stacks-geometry-lemma-deformation-category} と
Formal Deformation Theory,
Definition \ref{formal-defos-definition-deformation-category}，
および Lemmas \ref{formal-defos-lemma-RS-implies-S1-S2} と
\ref{formal-defos-lemma-versal-object-existence} から従う．
Formal Deformation Theory, Lemma
\ref{formal-defos-lemma-minimal-versal} の一意性の結果により，
「極小」半普遍変形環 $A$ が $\mathcal{X}$ の $x_0$ において存在し，
$\mathcal{X}$ の $x_0$ における他の任意の半普遍変形環は，
$A[[t_1, \ldots, t_r]]$ とある $r$ に対して同型である．
これは第二の主張を明らかに含意する．
\end{proof}

\begin{lemma}
\label{stacks-geometry-lemma-versal-ring-field-extension}
Situation \ref{stacks-geometry-situation-versal} において，
$x_0 : \Spec(k) \to \mathcal{X}$ を射とし，
$k$ は $S$ 上有限型な体とする．
$l/k$ を体の有限拡大とし，誘導される射を
% P12-ERR-0145: Japanese supplies the relation omitted by the frozen English notation clause.
$x_{l, 0} : \Spec(l) \to \mathcal{X}$ と書く．
$x_0$ における $\mathcal{X}$ の半普遍変形環 $A$ が与えられれば，
$x_{l, 0}$ における $\mathcal{X}$ の半普遍変形環 $A'$ であって，
与えられた体拡大 $l/k$ を誘導し，かつ
$\mathfrak m_{A'}$-進位相に関して形式的に平滑である
$S$-代数準同型 $A \to A'$ をもつものが存在する．
\end{lemma}

\begin{proof}
これは Artin's Axioms, Lemma \ref{artin-lemma-change-of-field}
および Formal Deformation Theory, Lemma
\ref{formal-defos-lemma-change-of-field-versal-ring}
から直ちに従う．
（Artin's Axioms, Lemma
\ref{artin-lemma-algebraic-stack-RS} により
$\mathcal{X}$ が (RS) を満たすことも用いる．）
\end{proof}

\begin{lemma}
\label{stacks-geometry-lemma-compare-versal-ring-completion}
Situation \ref{stacks-geometry-situation-versal} において，
$x : U \to \mathcal{X}$ を射とし，$U$ は $S$ 上局所有限型な
スキームとする．$u_0 \in U$ を有限型点とする．
$k = \kappa(u_0)$ とおき，誘導される射を
% P12-ERR-0146: Japanese supplies the relation omitted by the frozen English notation clause.
$x_0 : \Spec(k) \to \mathcal{X}$ と書く．次は同値である：
\begin{enumerate}
\item $x$ は $u_0$ において半普遍である
（Artin's Axioms, Definition \ref{artin-definition-versal}）；
\item $\hat x : \mathcal{F}_{U, k, u_0} \to
\mathcal{F}_{\mathcal{X}, k, x_0}$ は平滑である；
\item $x|_{\Spec(\mathcal{O}_{U, u_0}^\wedge)}$
に付随する形式対象は半普遍である；
\item
% P12-ERR-0147: The frozen source says a neighbourhood of the morphism x inside U; Japanese preserves x rather than silently substituting u_0.
$x$ の開近傍 $U' \subset U$ であって，
$x|_{U'} : U' \to \mathcal{X}$ が平滑となるものが存在する．
\end{enumerate}
さらに，この場合，完備化 $\mathcal{O}_{U, u_0}^\wedge$ は
$x_0$ における $\mathcal{X}$ の半普遍変形環である．
\end{lemma}

\begin{proof}
$U \to S$ は局所有限型であるから（そのような射の合成として），
$\Spec(k) \to S$ も有限型である（やはり合成として）．
したがって主張は意味をもつ．(1) と (2) の同値性は，
$x$ が $u_0$ において半普遍であることの定義である．
(1) と (3) の同値性は Artin's Axioms, Lemma
\ref{artin-lemma-versality-matches} である．
よって (1)，(2)，(3) は同値である．

\medskip\noindent
$x|_{U'}$ が平滑ならば，Artin's Axioms, Lemma
\ref{artin-lemma-formally-smooth-on-deformation-categories}
により関手
$\hat x : \mathcal{F}_{U, k, u_0} \to
\mathcal{F}_{\mathcal{X}, k, x_0}$ は平滑である．
したがって (4) は (1)，(2)，(3) を含意する．
逆に，$x$ が $u_0$ において半普遍であると仮定する．
$V$ がスキームであるような全射平滑射
$y : V \to \mathcal{X}$ を選ぶ．
$Z = V \times_\mathcal{X} U$ とおき，$u_0$ の上にある
有限型点 $z_0 \in |Z|$ を選ぶ
（これは Morphisms of Spaces, Lemma
\ref{spaces-morphisms-lemma-finite-type-points-surjective-morphism}
により可能である）．Artin's Axioms, Lemma
\ref{artin-lemma-base-change-versal} により，射 $Z \to V$ は
$z_0$ において平滑である．定義により，$z_0$ の開近傍
$W \subset Z$ であって $W \to V$ が平滑となるものを選べる．
$Z \to U$ は開写像なので，$U' \subset U$ を $W$ の像とする．
すると，スタックの平滑射の定義から
$U' \to \mathcal{X}$ は平滑である．

\medskip\noindent
最後の主張は，$\mathcal{O}_{U, u_0}^\wedge$ が
$\mathcal{F}_{U, k, u_0}$ をプロ表現することから，
定義により従う．
\end{proof}

\begin{lemma}
\label{stacks-geometry-lemma-characterize-smoothness}
% P12-ERR-0148: Japanese joins the frozen sentence-fragment opening grammatically.
Situation \ref{stacks-geometry-situation-versal} において，
$x_0 : \Spec(k) \to \mathcal{X}$ を射とし，
$\Spec(k) \to S$ は像を $s$ とする有限型射であるとする．
$A$ を $x_0$ における $\mathcal{X}$ の半普遍変形環とする．
次は同値である：
\begin{enumerate}
\item $x_0$ は $\mathcal{X} \to S$ の平滑軌跡に属する
（Morphisms of Stacks, Lemma
\ref{stacks-morphisms-lemma-where-smooth}）；
\item $\mathcal{O}_{S, s} \to A$ は
$\mathfrak m_A$-進位相に関して形式的に平滑である；
\item $\mathcal{F}_{\mathcal{X}, k, x_0}$ は障害をもたない．
\end{enumerate}
\end{lemma}

\begin{proof}
(2) と (3) の同値性は Formal Deformation Theory, Lemma
\ref{formal-defos-lemma-smooth-power-series-classical}
から直ちに従う．

\medskip\noindent
$\mathcal{O}_{S, s} \to A$ が
$\mathfrak m_A$-進位相に関して形式的に平滑であることと，
$\mathcal{O}_{S, s} \to A' = A[[t_1, \ldots, t_r]]$ が
$\mathfrak m_{A'}$-進位相に関して形式的に平滑であることは同値である．
したがって Lemma \ref{stacks-geometry-lemma-versal-ring} により，(2) は
半普遍変形環の選択に依存しない．次に，$l/k$ を有限拡大とし，
Lemma \ref{stacks-geometry-lemma-versal-ring-field-extension} のような
$A \to A'$ を選ぶ．$\mathcal{O}_{S, s} \to A$ が
$\mathfrak m_A$-進位相に関して形式的に平滑ならば，
More on Algebra, Lemma
\ref{more-algebra-lemma-compose-formally-smooth} により
$\mathcal{O}_{S, s} \to A'$ は
$\mathfrak m_{A'}$-進位相に関して形式的に平滑である．
逆に，$\mathcal{O}_{S, s} \to A'$ が
$\mathfrak m_{A'}$-進位相に関して形式的に平滑ならば，
$\mathcal{O}_{S, s}^\wedge \to A'$ と $A \to A'$ は正則であり
（More on Algebra, Proposition
\ref{more-algebra-proposition-fs-regular}），
したがって $\mathcal{O}_{S, s}^\wedge \to A$ は正則であり
（More on Algebra, Lemma
\ref{more-algebra-lemma-regular-permanence}），
それゆえ $\mathcal{O}_{S, s} \to A$ は
$\mathfrak m_A$-進位相に関して形式的に平滑である
（先と同じ補題）．よって，$k$ と $x_0$ に対して (2) と (1)
が同値であることと，$l$ と $x_{0, l}$ に対して同値であることは等価である．

\medskip\noindent
$U$ をスキームとし，$\Spec(k) \times_\mathcal{X} U$ が空でないような
平滑射 $U \to \mathcal{X}$ を選ぶ．有限拡大 $l/k$ と点
$w_0 : \Spec(l) \to \Spec(k) \times_\mathcal{X} U$ を選び，
$u_0 \in U$ を $w_0$ の像とする．上の議論を $l/k$ および
$l/\kappa(u_0)$ に適用することにより，$u_0$ の場合へ帰着できる．
したがって Lemma \ref{stacks-geometry-lemma-compare-versal-ring-completion} を用いて
$A = \mathcal{O}_{U, u_0}^\wedge$ と仮定してよい．
$x_0$ が $\mathcal{X} \to S$ の平滑軌跡に属することと，
$u_0$ が $U \to S$ の平滑軌跡に属することは同値である；
例えば Morphisms of Stacks, Lemma
\ref{stacks-morphisms-lemma-where-smooth} を参照せよ．
したがって (1) と (2) の同値性は More on Algebra, Lemma
\ref{more-algebra-lemma-formally-smooth-finite-type} から従う．
\end{proof}

\noindent
Artin 近似の一つの帰結を思い出しておく．

\begin{lemma}
\label{stacks-geometry-lemma-Artin-approximation-by-smooth-morphism}
% P12-ERR-0148: Japanese joins the second frozen sentence-fragment opening grammatically.
Situation \ref{stacks-geometry-situation-versal} において，
$x_0 : \Spec(k) \to \mathcal{X}$ を射とし，
$\Spec(k) \to S$ は像を $s$ とする有限型射であるとする．
$A$ を $x_0$ における $\mathcal{X}$ の半普遍変形環とする．
$\mathcal{O}_{S, s}$ が G-環ならば，始域がスキームである平滑射
$U \to \mathcal{X}$ と剰余体 $k$ をもつ点 $u_0 \in U$ であって，
\begin{enumerate}
\item $\Spec(k) \to U \to \mathcal{X}$ が与えられた射 $x_0$ と一致し，
\item 同型 $\mathcal{O}_{U, u_0}^\wedge \cong A$ が存在する
\end{enumerate}
ものを見いだせる．
\end{lemma}

\begin{proof}
$(\xi_n, f_n)$ を $A$ 上の半普遍形式対象とする．
Artin's Axioms, Lemma \ref{artin-lemma-effective} により，
$\xi = (A, \xi_n, f_n)$ は有効である．
仮定より $\mathcal{X}$ は $S$ 上局所有限表示であり
（Morphisms of Stacks, Lemma
\ref{stacks-morphisms-lemma-noetherian-finite-type-finite-presentation}
を用いる），したがって Limits of Stacks, Proposition
\ref{stacks-limits-proposition-characterize-locally-finite-presentation}
により極限を保つ．よって Artin's Axioms, Lemma
\ref{artin-lemma-approximate-versal} の Artin 近似から，
始域が有限型 $S$-スキームである射 $U \to \mathcal{X}$ と，
剰余体 $k$ をもつ点 $u_0 \in U$ で (1) と (2) を満たし，
さらに $U \to \mathcal{X}$ が $u_0$ において半普遍となるものを
見いだせる．Lemma \ref{stacks-geometry-lemma-compare-versal-ring-completion}
により，$U$ を縮小した後，$U \to \mathcal{X}$ が平滑であると
仮定してよい．
\end{proof}

\begin{remark}[半普遍変形環の強化]
\label{stacks-geometry-remark-upgrade}
Situation \ref{stacks-geometry-situation-versal} において，
$x_0 : \Spec(k) \to \mathcal{X}$ を射とし，
$k$ は $S$ 上有限型な体とする．
$A$ を $x_0$ における $\mathcal{X}$ の半普遍変形環とする．
Artin's Axioms, Lemma \ref{artin-lemma-effective} により，
半普遍形式対象は実際には $S$ 上の射
$$
\Spec(A) \longrightarrow \mathcal{X}
$$
から得られる．さらに，上の各結果はこの射と両立する形に強化できる．
% P12-ERR-0149: Japanese normalizes the frozen English word order.
その一覧は次のとおりである：
\begin{enumerate}
\item Lemma \ref{stacks-geometry-lemma-versal-ring} における同型
$A \cong A'[[t_1, \ldots, t_r]]$ または
$A' \cong A[[t_1, \ldots, t_r]]$ は，これらの射と両立するように選べる；
\item Lemma \ref{stacks-geometry-lemma-versal-ring-field-extension} における準同型
$A \to A'$ は，これらの射と両立するように選べる；
\item Lemma \ref{stacks-geometry-lemma-compare-versal-ring-completion} における射
$\Spec(\mathcal{O}_{U, u_0}^\wedge) \to \mathcal{X}$ は，標準射
$\Spec(\mathcal{O}_{U, u_0}^\wedge) \to U$ と与えられた射
$U \to \mathcal{X}$ との合成である；
\item Lemma \ref{stacks-geometry-lemma-Artin-approximation-by-smooth-morphism} における同型
$\mathcal{O}_{U, u_0}^\wedge \cong A$ は，
$\Spec(A) \to \mathcal{X}$ が前項の標準射に対応するように選べる．
\end{enumerate}
いずれの場合も，主張は写像が半普遍形式要素と両立することから従う．
ただし，暗黙に現れる図式が $2$-可換となるのは，
ある（非標準的な）$2$-射を選んだ後に限られることに注意する．
それでも，(1) における暗黙の写像 $A' \to A$ または $A \to A'$ は
形式的ホモトピーを除いて良定義である．
Formal Deformation Theory, Lemma
\ref{formal-defos-lemma-versal-unique-up-to-homotopy} を参照せよ．
\end{remark}

\begin{lemma}
\label{stacks-geometry-lemma-versal-ring-flat}
Situation \ref{stacks-geometry-situation-versal} において，
$x_0 : \Spec(k) \to \mathcal{X}$ を射とし，
$k$ は $S$ 上有限型な体とする．
$A$ を $x_0$ における $\mathcal{X}$ の半普遍変形環とする．
このとき Remark \ref{stacks-geometry-remark-upgrade} の射
$\Spec(A) \to \mathcal{X}$ は平坦である．
\end{lemma}

\begin{proof}
$S$ の像の点における局所環が G-環ならば，
これは Lemma \ref{stacks-geometry-lemma-Artin-approximation-by-smooth-morphism}
と，Noether 的局所環からその完備化への写像が平坦であることから
直ちに従う．一般の場合には，次のように示す．

\medskip\noindent
Step I. $A$ と $A'$ を $x_0$ における $\mathcal{X}$ の
二つの半普遍変形環とする．このとき，主張が $A$ に対して成り立つことと
$A'$ に対して成り立つことは同値である．実際，必要なら
$A$ と $A'$ を入れ替えることにより，合成
$$
\Spec(A') \to \Spec(A) \to \mathcal{X}
$$
が射 $\Spec(A') \to \mathcal{X}$ となるような形式的平滑射
$\varphi : A \to A'$ が存在すると仮定してよい．
Lemma \ref{stacks-geometry-lemma-versal-ring} および
Remark \ref{stacks-geometry-remark-upgrade} を参照せよ．
$A \to A'$ は忠実平坦なので，Morphisms of Stacks, Lemmas
\ref{stacks-morphisms-lemma-composition-flat} および
\ref{stacks-morphisms-lemma-flat-permanence} から同値性が得られる．

\medskip\noindent
Step II. $l/k$ を体の有限拡大とし，誘導される射を
$x_{l, 0} : \Spec(l) \to \mathcal{X}$ とする．
$A$ を $x_0$ における $\mathcal{X}$ の半普遍変形環とし，
$A \to A'$ を Lemma \ref{stacks-geometry-lemma-versal-ring-field-extension}
のように選ぶ．このときも，合成
$$
\Spec(A') \to \Spec(A) \to \mathcal{X}
$$
は射 $\Spec(A') \to \mathcal{X}$ である；
Remark \ref{stacks-geometry-remark-upgrade} を参照せよ．
先と同様に議論し，Step I により半普遍変形環の選択が無関係であることを
用いれば，補題が $x_0$ に対して成り立つことと
$x_{l, 0}$ に対して成り立つことは同値である．

\medskip\noindent
Step III. $U$ をスキームとし，全射平滑射
$U \to \mathcal{X}$ を選ぶ．すると
$Z = U \times_\mathcal{X} x_0$（空でない代数空間）上の
有限型点 $z_0$ を選べる．$u_0 \in U$ を $z_0$ の $U$ における像とする．
$z_0$ が閉点 $w_0 \in W$ の像となるようなスキームと全射 \'etale 射
$W \to Z$ を選ぶ
（Morphisms of Spaces, Section
\ref{spaces-morphisms-section-points-finite-type} を参照）．
$W \to \Spec(k)$ と $W \to U$ は有限型なので，
$\kappa(w_0)/k$ および $\kappa(w_0)/\kappa(u_0)$ は体の有限拡大である
（Morphisms, Section
\ref{morphisms-section-points-finite-type} を参照）．
Step II を二度適用して，$x_0$ を
$u_0 \to U \to \mathcal{X}$ で置き換えてよい．
すると，問題の射は合成
$$
\Spec(\mathcal{O}_{U, u_0}^\wedge) \to U \to \mathcal{X}
$$
である．
% P12-ERR-0150: Japanese normalizes the frozen English number disagreement.
第一の矢印は Noether 的局所環の完備化が平坦であること
（Algebra, Lemma \ref{algebra-lemma-completion-flat}）から平坦であり，
第二の矢印は平滑射が平坦であることから平坦である．
平坦性は合成で保たれるので，合成も平坦である．
\end{proof}

\begin{remark}
\label{stacks-geometry-remark-groupoid-defo}
Situation \ref{stacks-geometry-situation-versal} において，
$x_0 : \Spec(k) \to \mathcal{X}$ を射とし，
$k$ は $S$ 上有限型な体とする．
Lemma \ref{stacks-geometry-lemma-deformation-category} および
Formal Deformation Theory, Theorem
\ref{formal-defos-theorem-presentation-deformation-groupoid}
により，$\mathcal{F}_{\mathcal{X}, k, x_0}$ は
$\mathcal{C}_\Lambda$ 上の関手における平滑な
プロ表現可能亜群による表示をもつ．
定義を展開すると，次を選べるということである：
\begin{enumerate}
\item 剰余体 $k$ をもつ Noether 的完備局所 $\Lambda$-代数 $A$ と，
$A$ 上の $\mathcal{F}_{\mathcal{X}, k, x_0}$ の
半普遍形式対象 $\xi$；
\item 剰余体 $k$ をもつ Noether 的完備局所 $\Lambda$-代数 $B$ と同型
$$
\underline{B}|_{\mathcal{C}_\Lambda}
\longrightarrow
\underline{A}|_{\mathcal{C}_\Lambda}
\times_{\underline{\xi}, \mathcal{F}_{\mathcal{X}, k, x_0}, \underline{\xi}}
\underline{A}|_{\mathcal{C}_\Lambda}
$$
\end{enumerate}
二つの射影は形式的平滑写像
$t : A \to B$ および $s : A \to B$ に対応する
（$\xi$ が半普遍であるため）．
写像 $c : B \to B \widehat{\otimes}_{s, A, t} B$ が存在し，
$(A, B, s, t, c)$ を，剰余体 $k$ をもつ Noether 的完備局所
$\Lambda$-代数の圏における余亜群とする
（プロ表現可能関手上では，この写像は
Formal Deformation Theory, Lemma
\ref{formal-defos-lemma-presentation-construction} で構成される）．
最後に，引用した定理によれば，$\xi$ は
$$
[\underline{A}|_{\mathcal{C}_\Lambda} / \underline{B}|_{\mathcal{C}_\Lambda}]
\longrightarrow
\mathcal{F}_{\mathcal{X}, k, x_0}
$$
という $\mathcal{C}_\Lambda$ 上に余ファイバー化された亜群の同値を
誘導する．実際，完備化された圏
$\widehat{\mathcal{C}}_\Lambda$ 上に余ファイバー化された亜群の同値
$$
[\underline{A}/\underline{B}]
\longrightarrow
\widehat{\mathcal{F}}_{\mathcal{X}, k, x_0}
$$
も得られる
（これが成り立つ理由については Formal Deformation Theory, Section
\ref{formal-defos-section-prorepresentable-groupoids-in-functors}
の議論を参照）．もちろん $A$ は $x_0$ における
$\mathcal{X}$ の半普遍変形環である．
\end{remark}







% BEGIN EXPANDED MODULE ch107_02.tex
\section{代数スタックの成分の重複度}
\label{stacks-geometry-section-multiplicities}

\noindent
$X$ を局所 Noether スキームとすると，$X$ を（単なる位相空間とみなして）
既約成分の合併として，たとえば $X = \bigcup T_i$ と書くことができる．
各既約成分は一意な生成点 $\xi_i$ の閉包であり，局所環
$\mathcal O_{X,\xi_i}$ は Artin 局所環である．{\it $X$ の $T_i$ に沿う重複度}
または {\it $T_i$ の $X$ における重複度} を
$$
m_{T_i, X} = \text{長さ}_{\mathcal O_{X, \xi_i}} \mathcal O_{X, \xi_i}
$$
によって定義できる．言い換えれば，これはこの Artin 局所環の長さである．
次と比較されたい：
Chow Homology, Section \ref{chow-section-cycle-of-closed-subscheme}．

\medskip\noindent
ここでの目標は，この定義を局所 Noether 代数スタックへ一般化することである．
$\mathcal{X}$ がスタックならば，その位相空間 $|\mathcal{X}|$
（Properties of Stacks, Definition
\ref{stacks-properties-definition-topological-space} を参照）は
局所 Noether である
（Morphisms of Stacks, Lemma \ref{stacks-morphisms-lemma-Noetherian-topology}）．
$|\mathcal{X}|$ の既約成分を $\mathcal{X}$ の既約成分と
呼ぶこともある．
$\mathcal{X}$ が準分離ならば，$|\mathcal{X}|$ は
sober である（Morphisms of Stacks, Lemma
\ref{stacks-morphisms-lemma-sober-qs}）が，
非準分離の場合にはそうとは限らない．たとえば，非準分離な代数空間
$X = \mathbf{A}^1_\mathbf{C}/\mathbf{Z}$ を考えればよい．
さらに，$|\mathcal{X}|$ 上には，その茎を用いて重複度を
定義できるような構造層は存在しない．

\begin{lemma}
\label{stacks-geometry-lemma-map-of-components}
$f : U \to \mathcal{X}$ を，スキームから局所 Noether 代数スタックへの
滑らかな射とする．$|U|$ の任意の既約成分の像の閉包は
$|\mathcal{X}|$ の既約成分である．
$U \to \mathcal{X}$ が全射ならば，$|\mathcal{X}|$ のすべての既約成分が
このようにして得られる．
\end{lemma}

\begin{proof}
写像 $|U| \to |\mathcal{X}|$ は，Properties of Stacks, Lemma
\ref{stacks-properties-lemma-topology-points} により連続かつ開である．
$T \subset |U|$ を既約成分とする．$U$ は局所 Noether なので，
空でないアフィン開集合 $W \subset U$ で $T$ に含まれるものをとることができる．
このとき $f(T) \subset |\mathcal{X}|$ は既約であり，空でない開部分集合
$f(W)$ を含む．したがって $f(T)$ の閉包は既約であり，
空でない開集合を含む．ゆえに，この閉包は既約成分である．

\medskip\noindent
$U \to \mathcal{X}$ が全射であると仮定し，$Z \subset |\mathcal{X}|$ を
既約成分とする．Noether 開部分集合 $V$ を $|\mathcal{X}|$ の中にとり，
$Z$ と交わるようにする．$V$ から他の既約成分を除けば，
$V \subset Z$ と仮定してよい．空でない開集合
$f^{-1}(V) \subset |U|$ の既約成分を一つとり，$T \subset |U|$ をその閉包とする．
これは $|U|$ の既約成分であり，$f(T)$ の閉包は $Z$ と，
$T$ の選び方により一致しなければならない．
\end{proof}

\noindent
前の補題は，とくに局所 Noether スキーム間の滑らかな射の場合に
適用できる．この特別な場合は，次の補題の主張で
暗黙に用いられている．

\begin{lemma}
\label{stacks-geometry-lemma-multiplicities}
$U \to X$ を局所 Noether スキームの滑らかな射とする．
$T'$ を $U$ の既約成分とする．% P12-ERR-0151
$T$ を $X$ の既約成分で，$T'$ の像の閉包として得られるものとする．
このとき $m_{T', U} = m_{T, X}$ である．
\end{lemma}

\begin{proof}
$\xi'$ で $T'$ の生成点を，$\xi$ で
$T$ の生成点を表す．$A = \mathcal{O}_{X, \xi}$ および
$B = \mathcal{O}_{U, \xi'}$ とおく．示すべきことは
$\text{長さ}_A A = \text{長さ}_B B$ である．
$A \to B$ は平坦な局所環準同型
（滑らかな射は平坦だからである）なので，
$$
\text{長さ}_A(A) \text{長さ}_B(B/\mathfrak m_A B) =
\text{長さ}_B(B)
$$
が Algebra, Lemma \ref{algebra-lemma-pullback-module} により成り立つ．したがって，
$\mathfrak m_A B = \mathfrak m_B$，同値なこととして
$B/\mathfrak m_A B$ が被約であることを示せば十分である．
$U \to X$ は滑らかなので，その基底変換
$U_{\xi} \to \Spec \kappa(\xi)$ も滑らかである．$U_{\xi}$ は
体上の滑らかなスキームなので被約であり，したがって任意の点における
その局所環も被約である．% P12-ERR-0152
（Varieties, Lemma \ref{varieties-lemma-smooth-geometrically-normal}）．
とくに，
$$
B/\mathfrak m_A B =
\mathcal{O}_{U, \xi'}/\mathfrak m_{X, \xi}\mathcal{O}_{U, \xi'} =
\mathcal{O}_{U_\xi, \xi'}
$$
は被約であり，これが示すべきことであった．
\end{proof}

\noindent
この結果を用いると，滑らかな位相で局所的に見ることにより，
重複度の適切な概念が存在することを示せる．

\begin{lemma}
\label{stacks-geometry-lemma-multiplicity}
$U_1 \to \mathcal{X}$ および $U_2 \to \mathcal{X}$ を，スキームから
局所 Noether 代数スタック $\mathcal{X}$ への二つの滑らかな射とする．
$T_1'$ および $T_2'$ を，それぞれ $|U_1|$
および $|U_2|$ の既約成分とする．
$T_1'$ と $T_2'$ の像の閉包が，ともに同じ既約成分 $T$ であり，
これが $|\mathcal{X}|$ の既約成分であると仮定する．
このとき $m_{T_1', U_1} = m_{T_2', U_2}$ である．
\end{lemma}

\begin{proof}
$V_1$ および $V_2$ を，それぞれ $T_1'$ および $T'_2$ の稠密部分集合で，
それぞれ $U_1$ および $U_2$ において開なものとする
（Lemma \ref{stacks-geometry-lemma-map-of-components} の証明を参照）．
$|V_1|$ と $|V_2|$ の $|\mathcal{X}|$ における像は，
既約部分集合 $T$ の空でない開部分集合であり，したがって
その共通部分は空でない．
Properties of Stacks, Lemma \ref{stacks-properties-lemma-points-cartesian} により，
写像 $|V_1 \times_\mathcal{X} V_2| \to |V_1| \times_{|\mathcal{X}|} |V_2|$
は全射である．したがって $V_1 \times_\mathcal{X} V_2$ は
空でない代数空間であるから，始域が（空でない）スキームである
\'etale 全射 $V \to V_1 \times_\mathcal{X} V_2$ をとることができる．
$T'$ を $V$ の任意の既約成分とすると，
Lemma \ref{stacks-geometry-lemma-map-of-components} により，$T'$ の $U_1$
（それぞれ $U_2$）における像の閉包は $T'_1$
（それぞれ $T'_2$）に等しい．

\medskip\noindent
Lemma \ref{stacks-geometry-lemma-multiplicities} を二度適用すると，
$$
m_{T_1', U_1} = m_{T', V} = m_{T_2', U_2},
$$
を得る．これが示すべきことであった．
\end{proof}

\noindent
以上で，次の定義が意味をもつことを示すために
十分な準備が整った．

\begin{definition}
\label{stacks-geometry-definition-multiplicity}
$\mathcal{X}$ を局所 Noether 代数スタックとする．
$T \subset |\mathcal{X}|$ を既約成分とする．
{\it $T$ の $\mathcal{X}$ における重複度} を
$m_{T, \mathcal{X}} = m_{T', U}$ により定義する．ここで $f : U \to \mathcal{X}$ は
スキームからの滑らかな射であり，$T' \subset |U|$ は
$f(T') \subset T$ を満たす既約成分である．
\end{definition}

\noindent
これは，$f : U \to \mathcal{X}$ の選び方にも，既約成分 $T'$ で
$T$ へ写るものの選び方にも依存しない．これは Lemmas
\ref{stacks-geometry-lemma-map-of-components} および \ref{stacks-geometry-lemma-multiplicity} による．

\medskip\noindent
最後に，$\mathcal{X}$ の既約成分を閉部分スタックとみなすと
便利なことがあると注意しておく．
そのため，$T \subset |\mathcal{X}|$ が既約成分ならば，
一意な被約閉部分スタック $\mathcal{T} \subset \mathcal{X}$ で
$|\mathcal{T}| = T$ を満たすものを考えることができる．Properties of Stacks, Definition
\ref{stacks-properties-definition-reduced-induced-stack} を参照されたい．
$\mathcal{X}$ が準分離ならば，
既約成分は整スタックである．詳しくは
Morphisms of Stacks, Section \ref{stacks-morphisms-section-integral-stacks}
を参照されたい．



\section{形式的分枝と重複度}
\label{stacks-geometry-section-formal-branches}

\noindent
Definition \ref{stacks-geometry-definition-multiplicity} で与えた既約成分の重複度の概念と，
有限型の点における $\mathcal{X}$ の versal 環（のスペクトル）の
既約成分の重複度という関連する概念との比較を用意しておくと便利である．

\medskip\noindent
Situation \ref{stacks-geometry-situation-versal} において，$x_0 : \Spec(k) \to \mathcal{X}$ を
射とし，$k$ は $S$ 上有限型の体であるとする．
$A$, $A'$ を $\mathcal{X}$ の $x_0$ における versal 環とする．
必要なら $A$ と $A'$ を入れ替えることにより，versal 形式対象と両立する
形式的に滑らかな\footnote{
$A'$ が $A$ 上の冪級数環と同型になるという意味である．}
写像 $\varphi : A \to A'$ が存在することがわかる．Lemma
\ref{stacks-geometry-lemma-versal-ring} および
Remark \ref{stacks-geometry-remark-upgrade} を参照されたい．
さらに，$\varphi$ は形式的ホモトピーを除いてよく定まる．
Formal Deformation Theory, Lemma
\ref{formal-defos-lemma-versal-unique-up-to-homotopy} を参照されたい．
とくに，$\varphi(\mathfrak p)A'$ は
Formal Deformation Theory, Lemma
\ref{formal-defos-lemma-homotopic-minimal-prime} により，
$A'$ のよく定まったイデアルである．
$A \to A'$ は形式的に滑らかなので，実際
$\varphi(\mathfrak p)A'$ は $A'$ の極小素イデアルであり，
$A'$ のすべての極小素イデアルは，一意な極小素イデアル
$\mathfrak p \subset A$ に対してこの形になる（これらはすべて，
$A'$ を $A$ 上の冪級数環として書けば容易に示せる）．
したがって，極小素イデアルが既約成分に対応することを思い出せば，
次の定義は意味をもつ．

\begin{definition}
\label{stacks-geometry-definition-formal-branches}
$\mathcal{X}$ を，局所 Noether スキーム $S$ 上局所有限型の
代数スタックとする．$x_0 : \Spec(k) \to \mathcal{X}$ を射とし，
$k$ は $S$ 上有限型の体であるとする．% P12-ERR-0153
{\it $\mathcal{X}$ の $x_0$ を通る形式的分枝} とは，
$\Spec(A)$ の既約成分全体の集合である．これは $\mathcal{X}$ の
$x_0$ における versal 環を任意に選んで得られ，
上で述べた手続きによって異なる $A$ の選択の間で同一視する．% P12-ERR-0154
\end{definition}

\noindent
Definition \ref{stacks-geometry-definition-formal-branches} の状況で有限次拡大
$l/k$ が与えられたとする．
$x_{l, 0} : \Spec(l) \to \mathcal{X}$ を，
$\Spec(l) \to \Spec(k)$ と $x_0$ の合成とする．
$A \to A'$ を Lemma \ref{stacks-geometry-lemma-versal-ring-field-extension} におけるものとする．
$A \to A'$ は忠実平坦なので，射
$$
\Spec(A') \to \Spec(A)
$$
は，既約成分（の生成点）を
既約成分（の生成点）へ送る．
これは全射となるが，
一般には全単射とはならない．
言い換えると，全射
$$
\text{形式的分枝：}\mathcal{X}\text{，通過点 }x_{l, 0}
\longrightarrow
\text{形式的分枝：}\mathcal{X}\text{，通過点 }x_0
$$
を得る．$l/k$ が純非分離ならば，
この写像は単射でもあることがわかる（これが必要になった場合には，
ここに正確な主張と証明を追加する）．

\begin{lemma}
\label{stacks-geometry-lemma-branches}
Definition \ref{stacks-geometry-definition-formal-branches} の状況では，$\mathcal{X}$ の
$x_0$ を通る形式的分枝全体の集合から，$|\mathcal{X}|$ において
$x_0$ を含む $|\mathcal{X}|$ の既約成分全体の集合への標準的な全射が存在する．
\end{lemma}

\begin{proof}
$A$ を Definition \ref{stacks-geometry-definition-formal-branches} におけるものとし，
$\Spec(A) \to \mathcal{X}$ を Remark \ref{stacks-geometry-remark-upgrade} におけるものとする．
$\Spec(A)$ の既約成分の生成点は，
$|\mathcal{X}|$ の既約成分の生成点へ写ると主張する．
スキーム $U$ と全射な滑らかな射
$U \to \mathcal{X}$ をとる．次の図式を考える：
$$
\xymatrix{
\Spec(A) \times_\mathcal{X} U \ar[d]_p \ar[r]_-q & U \ar[d]^f \\
\Spec(A) \ar[r]^j & \mathcal{X}
}
$$
Lemma \ref{stacks-geometry-lemma-versal-ring-flat} により，$j$ は平坦である．
したがって $q$ は平坦である．一方，$f$ は全射かつ滑らかなので，
$p$ は全射かつ滑らかである．これより，既約成分の任意の生成点
$\eta \in \Spec(A)$ は，余次元 $0$ の点 $\eta'$ で，代数空間
$\Spec(A) \times_\mathcal{X} U$ の点であるものの像である（記法については
Properties of Spaces, Section \ref{spaces-properties-section-generic-points}
を参照し，\'etale 局所環に下降定理を用いる）．
$q$ は平坦なので，$q(\eta')$ は余次元 $0$ の $U$ の点である
（同じ議論による）．
$U$ はスキームなので，$q(\eta')$ は
$U$ の既約成分の生成点である．したがって
$q(\eta')$ の $|\mathcal{X}|$ における像の閉包は，
Lemma \ref{stacks-geometry-lemma-map-of-components} により既約成分である．これで主張が示された．

\medskip\noindent
この主張が，求める写像を定義する方法を与えることは明らかである．
その全射性を示すため，$u_0 \in U$ を，$x_0$ へ
$|\mathcal{X}|$ において写るようにとる．
アフィン開集合 $U' \subset U$ で $u_0$ の近傍であるものをとる．
$U'$ を縮小すれば，$U'$ のすべての既約成分が
$u_0$ を通ると仮定してよい．次に $\mathcal{X}$ を，
$|U'| \to |\mathcal{X}|$ の像に対応する開部分スタックで置き換えてよい．
したがって $U$ はアフィンで，点 $u_0$ をもち，これは
$x_0 \in |\mathcal{X}|$ へ写り，$U$ のすべての既約成分が
$u_0$ を通ると仮定してよい．% P12-ERR-0155
Properties of Stacks, Lemma
\ref{stacks-properties-lemma-points-cartesian} により，
点 $t \in |\Spec(A) \times_\mathcal{X} U|$ が存在し，
$\Spec(A)$ の閉点および $u_0$ へ写る．
平坦な局所環準同型に下降定理を用いると，
$$
A \longrightarrow
\mathcal{O}_{\Spec(A) \times_\mathcal{X} U, \overline{t}}
\longleftarrow \mathcal{O}_{U, u_0}
$$
$\mathcal{O}_{U, u_0}$ のすべての極小素イデアルは，
中央の局所環の極小素イデアルの像であり，そのような極小素イデアルは
$A$ の極小素イデアルへ写ることがわかる．
これで全射性が示された．いくつかの詳細は省略する．
\end{proof}

\noindent
$A$ を Noether 完備局所環とする．このとき $\Spec(A)$ の
既約成分には重複度がある．Section \ref{stacks-geometry-section-multiplicities} の
導入部を参照されたい．
$A' = A[[t_1, \ldots, t_r]]$ ならば，射
$\Spec(A') \to \Spec(A)$ は，重複度を保つ既約成分間の
全単射を誘導する（容易な証明は省略する）．
このことと，Definition \ref{stacks-geometry-definition-formal-branches} に先立つ議論により，
次の定義は意味をもつ．

\begin{definition}
\label{stacks-geometry-definition-multiplicity-formal-branches}
$\mathcal{X}$ を，局所 Noether スキーム $S$ 上局所有限型の
代数スタックとする．$x_0 : \Spec(k) \to \mathcal{X}$ を射とし，
$k$ は $S$ 上有限型の体であるとする．% P12-ERR-0156
{\it $\mathcal{X}$ の $x_0$ を通る形式的分枝の重複度} とは，
$\Spec(A)$ の対応する既約成分の重複度である．ここで versal 環は
$\mathcal{X}$ の $x_0$ におけるものを任意に選ぶ
（上の議論を参照）．
\end{definition}

\begin{lemma}
\label{stacks-geometry-lemma-branches-multiplicity}
$\mathcal{X}$ を，局所 Noether スキーム $S$ 上局所有限型の
代数スタックとする．$x_0 : \Spec(k) \to \mathcal{X}$ を射とし，
$k$ は $S$ 上有限型の体で，像が $s \in S$ であるとする．% P12-ERR-0157
$\mathcal{O}_{S, s}$ が G-環ならば，
Lemma \ref{stacks-geometry-lemma-branches} の写像は重複度を保つ．
\end{lemma}

\begin{proof}
Lemma \ref{stacks-geometry-lemma-Artin-approximation-by-smooth-morphism} により，
滑らかな射 $U \to \mathcal{X}$ で，$U$ がスキームであるものと，
$k$-値点 $u_0$ で，$U$ の点であり，
$\mathcal{O}_{U, u_0}^\wedge$ が $\mathcal{X}$ の
$x_0$ における versal 環となるものが存在すると仮定してよい．
Lemma \ref{stacks-geometry-lemma-branches} の証明における写像の構成
（$A = \mathcal{O}_{U, u_0}^\wedge$ なので大幅に簡単になる）により，
次を示せば十分である：$U$ の既約成分で $u_0$ を通るものの
重複度は，そこへ写る
$\Spec(\mathcal{O}_{U, u_0}^\wedge)$ の任意の既約成分の
重複度と等しい．

\medskip\noindent
可換環論の言葉に翻訳すると，次のようになる：
$C = \mathcal{O}_{U, u_0}$ とおく．これは
$\mathcal{O}_{S, s}$ 上本質的に有限型であり，したがって G-環である
（More on Algebra, Proposition
\ref{more-algebra-proposition-finite-type-over-G-ring}）．
$A = C^\wedge$ とおく．したがって $C \to A$ は正則環準同型である．
$\mathfrak q \subset C$ を極小素イデアルとし，極小素イデアル
$\mathfrak p \subset A$ で $\mathfrak q$ 上にあるものをとる．このとき
% P12-ERR-0158
$$
R = C_\mathfrak p \longrightarrow A_\mathfrak p = R'
$$
は Artin 局所環の正則環準同型である．このような環準同型については，
常に
$$
\text{長さ}_R R = \text{長さ}_{R'} R'
$$
が成り立つ．左辺は $U$ 上の当該成分の重複度であり，
右辺は $\Spec(A)$ 上の当該成分の重複度なので，
これが示すべきことである．
この等式を見るため，まず
$$
\text{長さ}_R(R) \text{長さ}_{R'}(R'/\mathfrak m_R R') =
\text{長さ}_{R'}(R')
$$
を Algebra, Lemma \ref{algebra-lemma-pullback-module} により用いる．したがって，
$\mathfrak m_R R' = \mathfrak m_{R'}$ を示せば十分であり，これは
零次元局所環の正則準同型であることから従う．
\end{proof}
% END EXPANDED MODULE ch107_02.tex
% BEGIN EXPANDED MODULE ch107_03.tex
\section{代数スタックの次元論}
\label{stacks-geometry-section-dimension-of-algebraic-stacks}

\noindent
% P12-ERR-0159: Japanese resolves the frozen relative-clause antecedent in the attribution sentence.
代数スタックの次元論について，われわれが把握している文献上の主要な結果は
\cite{Osserman} によるものであり，そこでは余次元と相対次元の概念が
研究されている．ここでは，点における代数スタックの次元という概念を
より詳しく考察し，射の始域の点におけるファイバーの次元を
その始域および終域の次元と関係づけるさまざまな結果を証明する．また，
（下の Lemma \ref{stacks-geometry-lemma-dimension-formula}）の結果も証明する．これにより
（適切な仮定の下で）点における代数スタックの次元を
versal 環を用いて計算できる．

\medskip\noindent
われわれは結果を常に最適な形にすることを目指したわけではないが，
不必要な仮定を置くことは概ね避けてきた．しかし，代数スタックのある性質と，
そのスタックのある点における versal 環の性質とを比較するいくつかの結果では，
すべての局所環が $G$-環である局所 Noether 的スキームを基底とする
局所有限表示な代数スタックの場合に考察を限定した．これにより，
% P12-ERR-0160: Japanese uses singular agreement for the frozen result clause.
versal 環の幾何とスタック自体の幾何を比較する際に
Artin 近似を利用できるという利便性が得られる．しかし，この制限的な仮定は，
われわれが証明するさまざまな主張のすべてが成り立つために
必要ではないかもしれない．とはいえ，念頭に置いている応用では
この仮定が満たされるので，有用な場合にはこれを課すことにしている．

\medskip\noindent
$X$ がスキームであるとき，$\dim(X)$ を $X$ の次元として，
$X$ の台位相空間の Krull 次元と定義する．一方，$x$ が $X$ の点ならば，
$\dim_x (X)$ を $X$ の $x$ における次元として，$U$ が $X$ の開部分集合で
$x$ を含むときのその次元の最小値として定義する；
Properties, Definition \ref{properties-definition-dimension} を参照せよ．
関係式 $\dim(X) = \sup_{x \in X} \dim_x(X)$ が成り立つ；
Properties, Lemma \ref{properties-lemma-dimension} を参照せよ．
$X$ が局所 Noether 的ならば，$\dim_x(X)$ は，$x$ を通る $X$ の既約成分の
$x$ における次元の上限に一致する．

\medskip\noindent
$X$ が代数空間で $x \in |X|$ であるとする．このとき，
$\dim_x X = \dim_u U,$ と定義する．ここで $U$ は \'etale 全射
$U \to X$ をもつ任意のスキームであり，$u\in U$ は $x$ の上にある任意の点である；
Properties of Spaces, Definition
\ref{spaces-properties-definition-dimension-at-point} を参照せよ．
$\dim(X) = \sup_{x \in |X|} \dim_x(X)$ とおく；
Properties of Spaces, Definition \ref{spaces-properties-definition-dimension} を参照せよ．

\begin{remark}
\label{stacks-geometry-remark-dimension-algebraic-space}
一般には，代数空間 $X$ の点 $x$ における次元は，
台位相空間 $|X|$ の点 $x$ における次元と一致するとは限らない．
たとえば\ $k$ が標数零の体で $X =  \mathbf{A}^1_k / \mathbf{Z}$ ならば，
$X$ は各点で次元 $1$（$\mathbf{A}^1_k$ の次元）をもつが，
$|X|$ は密着位相をもち，したがって Krull 次元は零である．一方，
Algebraic Spaces, Example \ref{spaces-example-infinite-product} には，
各点では次元 $0$ であるにもかかわらず，$|X|$ が Krull 次元 $1$ の既約空間であり，
かつ一般点をもつ（したがって，任意の点における $|X|$ の次元が $1$ である）
代数空間の例が与えられている；この例についての議論は
Properties of Spaces, Section \ref{spaces-properties-section-dimension}
も参照せよ．

\medskip\noindent
他方，$X$ が Decent Spaces, Definition
\ref{decent-spaces-definition-very-reasonable} の意味で
{\it 良好な}代数空間ならば（特に，\ $X$ が準分離的ならば；
Decent Spaces, Section \ref{decent-spaces-section-reasonable-decent} を参照せよ），
実際，$X$ の $x$ における次元は $|X|$ の $x$ における次元と一致する；
Decent Spaces, Lemma \ref{decent-spaces-lemma-dimension-decent-space} を参照せよ．
\end{remark}

\noindent
代数スタックの次元を定義するためには，まず，始域が代数空間で
終域が代数スタックである射について，始域の点における相対次元という
概念を用意しておくと便利である．定義がやや込み入っているのは，
（スキームの場合とは異なり）代数スタックや代数空間の点を
体のスペクトルからの射として記述することはできず，
そのような射の同値類としてしか記述できないためにすぎない．

\begin{definition}
\label{stacks-geometry-definition-relative-dimension}
$f : T \to \mathcal{X}$ を，代数空間から代数スタックへの
局所有限型射とし，$t \in |T|$ を，像が $x \in | \mathcal{X}|$ である点とする．
このとき，$f$ の $t$ における {\it 相対次元}を
$\dim_t(T_x),$ と書き，次のように定義する：
体のスペクトルを始域とする射 $\Spec k \to \mathcal{X}$ を選び，
これが $x$ を表すようにする．さらに点
$t' \in |T \times_{\mathcal{X}} \Spec k|$ を選び，これが $t$ に写るようにする．
ここで用いる射影の終域は $|T|$ である
（そのような点 $t'$ は Properties of Stacks, Lemma
\ref{stacks-properties-lemma-points-cartesian} により存在する）；このとき
$$
\dim_t(T_x) = \dim_{t'}(T \times_{\mathcal{X}} \Spec k ).
$$
\end{definition}

\noindent
$T$ は代数空間で $\mathcal{X}$ は代数スタックなので，ファイバー積
$T \times_{\mathcal{X}} \Spec k$ は代数空間であることに注意せよ．したがって，
この定義の右辺の量は実際に定義されている（上の議論を参照せよ）．

\begin{remark}
\label{stacks-geometry-remark-relative-dimension}
(1)
（たとえば，スキームの局所有限型射の相対次元が基底変換で不変であることを
用いれば；たとえば Morphisms, Lemma
\ref{morphisms-lemma-dimension-fibre-after-base-change} を参照せよ）
$\dim_t(T_x)$ が，計算に用いた選択に依存せず良定義であることは容易に確かめられる．

\medskip\noindent
(2)
$\mathcal{X}$ も代数空間である場合，この定義が
Morphisms of Spaces, Definition
\ref{spaces-morphisms-definition-dimension-fibre} で与えられた
相対次元の定義と一致することは直ちに確かめられる．
\end{remark}

\noindent
次に，局所 Noether 的代数スタックの次元に関するわれわれの研究の
基礎となる次の補題を想起する．

\begin{lemma}
\label{stacks-geometry-lemma-behaviour-of-dimensions-wrt-smooth-morphisms}
$f: U \to X$ を局所 Noether 的代数空間の間の滑らかな射とし，
$u \in |U|$ の像を $x \in |X|$ とする．このとき
$$
\dim_u (U) = \dim_x(X) + \dim_{u} (U_x)
$$
が成り立つ．ここで $\dim_u (U_x)$ は
Definition \ref{stacks-geometry-definition-relative-dimension} によって定義される．
\end{lemma}

\begin{proof}
Morphisms of Spaces, Lemma
\ref{spaces-morphisms-lemma-smoothness-dimension-spaces} を参照せよ．
ここで用いた $\dim_u (U_x)$ の定義が，そこで用いられた定義と一致することは
Remark \ref{stacks-geometry-remark-relative-dimension} (2) による．
\end{proof}

\begin{lemma}
\label{stacks-geometry-lemma-dimension-for-stacks}
% P12-ERR-0162: Japanese completes the frozen conditional opening.
$\mathcal{X}$ を局所 Noether 的代数スタックとし，
$x \in |\mathcal{X}|$ とする．$U \to \mathcal{X}$ を代数空間から
$\mathcal{X}$ への滑らかな射とし，$u$ を $|U|$ の $x$ に写る任意の点とする．
このとき
$$
\dim_x(\mathcal{X}) =  \dim_u(U) - \dim_{u}(U_x)
$$
が成り立つ．ここで，相対次元 $\dim_u(U_x)$ は
Definition \ref{stacks-geometry-definition-relative-dimension} により定義され，
$x$ における $\mathcal{X}$ の次元は
Properties of Stacks, Definition
\ref{stacks-properties-definition-dimension-at-point} におけるものである．
\end{lemma}

\begin{proof}
Lemma \ref{stacks-geometry-lemma-behaviour-of-dimensions-wrt-smooth-morphisms} を用いると，右辺
$\dim_u(U) + \dim_u(U_x)$ が，滑らかな射 $U \to \mathcal{X}$ および
$u \in |U|$ の選択に依存しないことを確かめられる．% P12-ERR-0163
詳細は省略する．特に，$U$ はスキームであると仮定してよい．
この場合，$x$ の代表として合成射
$\Spec \kappa(u) \to U \to \mathcal{X}$ を選ぶことにより，
$\dim_u(U_x)$ を計算できる．ここで最初の射は，像が $u \in U$ である標準射である．
$R = U \times_{\mathcal{X}} U$ と書き，対角射を $e : U \to R$ と書くと，
基底変換による相対次元の不変性から
$\dim_u(U_x) = \dim_{e(u)}(R_u)$ が従う．したがって，右辺は
$\dim_u (U) - \dim_{e(u)}(R_u) = \dim_x(\mathcal{X})$ に等しく，所望の結果を得る．
\end{proof}

\begin{remark}
\label{stacks-geometry-remark-dimension-DM}
適切な意味で良好な（たとえば準分離的な）Deligne--Mumford スタックについては，
$\dim_x(\mathcal{X})$ は，再び位相的に定義された量
$\dim_x |\mathcal{X}|$ と一致する．しかし，より一般の Artin スタックでは，
通常これは成り立たない．たとえば
$\mathcal{X} = [\mathbf{A}^1/\mathbf{G}_m]$
（ある体上で，商は $\mathbf{G}_m$ の $\mathbf{A}^1$ への通常の乗法作用に関して取る）
とする．このとき $|\mathcal{X}|$ は二つの点からなり，一方は他方の特殊化である
（これらは $\mathbf{G}_m$ の $\mathbf{A}^1$ 上の二つの軌道に対応する）．したがって，
位相空間としての次元は $1$ であるが，両方の点
$x \in |\mathcal{X}|$ において $\dim_x (\mathcal{X}) = 0$ である．
（さらに極端な例は分類空間 $[\Spec k/\mathbf{G}_m]$ であり，
その唯一の点における次元は $-1$ に等しい．）
\end{remark}

\noindent
これで，Definition \ref{stacks-geometry-definition-relative-dimension} を
（局所 Noether 的な）代数スタックの間の（局所有限型）射の場合へ
拡張できる．

\begin{definition}
\label{stacks-geometry-definition-relative-dimension-for-stacks}
% P12-ERR-0165: Japanese uses the standard local-finite-type construction.
$f : \mathcal{T} \to \mathcal{X}$ を，局所 Noether 的代数スタックの間の
局所有限型射とし，$t \in |\mathcal{T}|$ を，像が
$x \in |\mathcal{X}|$ である点とする．このとき，$f$ の $t$ における
{\it 相対次元}を $\dim_t(\mathcal{T}_x),$ と書き，次のように定義する：
体のスペクトルを始域とする射 $\Spec k \to \mathcal{X}$ を選び，
これが $x$ を表すようにする．さらに点
$t' \in |\mathcal{T} \times_{\mathcal{X}} \Spec k|$ を選び，これが $t$ に写るようにする．
ここで用いる射影の終域は $|\mathcal{T}|$ である
% P12-ERR-0164: Japanese closes the frozen parenthetical citation before continuing.
（そのような点 $t'$ は Properties of Stacks, Lemma
\ref{stacks-properties-lemma-points-cartesian} により存在する；このとき
$$
\dim_t(\mathcal{T}_x) = \dim_{t'}(\mathcal{T} \times_{\mathcal{X}} \Spec k ).
$$
\end{definition}

\noindent
$\mathcal{T}$ と $\mathcal{X}$ は代数スタックなので，ファイバー積
$\mathcal{T}\times_{\mathcal{X}} \Spec k$ は代数スタックであり，
Morphisms of Stacks, Lemma
\ref{stacks-morphisms-lemma-locally-finite-type-locally-noetherian}
により局所 Noether 的である．したがって，この定義の右辺の量は
Properties of Stacks, Definition
\ref{stacks-properties-definition-dimension-at-point} により定義される．

\begin{remark}
\label{stacks-geometry-remark-dimension-tangent-space-well-defined}
標準的な操作により，$\dim_t(\mathcal{T}_x)$ が，計算に用いた選択に依存せず
良定義であることが分かる．
\end{remark}

\noindent
ここで，相対次元のいくつかの基本的な性質を確立する．これらは，
スキームの射の場合に対応する主張の明らかな一般化である．

\begin{lemma}
\label{stacks-geometry-lemma-base-change-invariance-of-relative-dimension}
局所 Noether 的スタックの射からなる Cartesian 図式
$$
\xymatrix{
\mathcal{T}' \ar[d]\ar[r] & \mathcal{T} \ar[d] \\
\mathcal{X}' \ar[r] & \mathcal{X}
}
$$
が与えられ，その垂直射は局所有限型であるとする．
$t' \in |\mathcal{T}'|$ とし，その像を $t$, $x'$, および $x$ と書く．
これらはそれぞれ $|\mathcal{T}|$, $|\mathcal{X}'|$, および $|\mathcal{X}|$ の点である．
このとき，$\dim_{t'}(\mathcal{T}'_{x'}) = \dim_{t}(\mathcal{T}_x).$
\end{lemma}

\begin{proof}
両辺は（定義により）同じファイバー積の次元として計算できる．
\end{proof}

\begin{lemma}
\label{stacks-geometry-lemma-behaviour-of-dimensions-wrt-smooth-morphisms-stacky}
$f: \mathcal{U} \to \mathcal{X}$ を局所 Noether 的代数スタックの間の
滑らかな射とし，$u \in |\mathcal{U}|$ の像を $x \in |\mathcal{X}|$ とする．
このとき
$$
\dim_u (\mathcal{U}) = \dim_x(\mathcal{X}) + \dim_{u} (\mathcal{U}_x).
$$
が成り立つ．
\end{lemma}

\begin{proof}
始域がスキームである滑らかな全射 $V \to \mathcal{U}$ を選び，
$v\in |V|$ を $u$ に写る点とする．このとき合成射
$V \to \mathcal{U} \to \mathcal{X}$ も滑らかであり，
Lemma \ref{stacks-geometry-lemma-behaviour-of-dimensions-wrt-smooth-morphisms} により
$\dim_x(\mathcal{X}) = \dim_v(V) - \dim_v(V_x),$
である一方，$\dim_u(\mathcal{U}) = \dim_v(V) - \dim_v(V_u).$
である．したがって
$$
\dim_u(\mathcal{U}) - \dim_x(\mathcal{X}) = \dim_v (V_x) - \dim_v (V_u).
$$

\medskip\noindent
代表 $\Spec k \to \mathcal{X}$ を選んで $x$ を表し，点
$v' \in | V \times_{\mathcal{X}} \Spec k|$ を $v$ の上に選ぶ．
その像を $u'$ と書くと，これは $|\mathcal{U}\times_{\mathcal{X}} \Spec k|$ の点であり，定義により
$\dim_u(\mathcal{U}_x) = \dim_{u'}(\mathcal{U}\times_{\mathcal{X}} \Spec k),$
かつ
$\dim_v(V_x) = \dim_{v'}(V\times_{\mathcal{X}} \Spec k).$

\medskip\noindent
ここで $V\times_{\mathcal{X}} \Spec k \to \mathcal{U}\times_{\mathcal{X}}\Spec k$
は滑らかな全射（そのような射の基底変換だからである）であり，
その始域は代数空間である（$V$ と $\Spec k$ はスキームであり，
$\mathcal{X}$ は代数スタックだからである）．したがって，再び定義により
\begin{align*}
\dim_{u'}(\mathcal{U}\times_{\mathcal{X}} \Spec k)
& =
\dim_{v'}(V\times_{\mathcal{X}} \Spec k) -
\dim_{v'}(V \times_{\mathcal{X}} \Spec k)_{u'}) \\
& = \dim_v(V_x) -
\dim_{v'}( (V\times_{\mathcal{X}} \Spec k)_{u'}).
\end{align*}
さて，$V\times_{\mathcal{X}} \Spec k \cong
V\times_{\mathcal{U}} (\mathcal{U}\times_{\mathcal{X}} \Spec k),$
なので，Lemma \ref{stacks-geometry-lemma-base-change-invariance-of-relative-dimension} から
$\dim_{v'}((V\times_{\mathcal{X}} \Spec k)_{u'})  = \dim_v(V_u).$
が従う．以上をすべて合わせると，
$$
\dim_u(\mathcal{U}) - \dim_x(\mathcal{X}) =
\dim_u(\mathcal{U}_x),
$$
を得る．これは求める等式である．
\end{proof}

\begin{lemma}
\label{stacks-geometry-lemma-relative-dimension-is-semi-continuous}
% P12-ERR-0161: Japanese uses the standard local-finite-type construction.
$f: \mathcal{T} \to \mathcal{X}$ を代数スタックの局所有限型射とする．
\begin{enumerate}
\item
関数 $t \mapsto \dim_t(\mathcal{T}_{f(t)})$ は $|\mathcal{T}|$ 上で上半連続である．
\item $f$ が滑らかならば，関数
$t \mapsto \dim_t(\mathcal{T}_{f(t)})$ は $|\mathcal{T}|$ 上で局所定数である．
\end{enumerate}
\end{lemma}

\begin{proof}
まず $\mathcal{T}$ がスキーム $T$ であると仮定する．始域がスキームである
滑らかな全射 $U \to \mathcal{X}$ を取り，$T' = T \times_{\mathcal{X}} U$ とおく．
$f': T' \to U$ を $f$ の $U$ 上での引き戻しとし，
$g: T' \to T$ を射影とする．

\medskip\noindent
Lemma \ref{stacks-geometry-lemma-base-change-invariance-of-relative-dimension} により，
$\dim_{t'}(T'_{f'(t')}) = \dim_{g(t')}(T_{f(g(t'))}),$
が $t' \in T'$ に対して成り立つ．一方，$g$ は滑らかな全射の基底変換なので
滑らかかつ全射であり，
$g$ が台位相空間上に誘導する写像は
Properties of Spaces, Lemma \ref{spaces-properties-lemma-topology-points}
により連続かつ開であり，しかも全射である．したがって，(1) が射 $f'$ に対して
% P12-ERR-0167: Japanese supplies a linking phrase at the frozen citation transition.
成り立つことは Morphisms of Spaces, Lemma
\ref{spaces-morphisms-lemma-openness-bounded-dimension-fibres} から従い，
(2) は Morphisms, Lemma
\ref{morphisms-lemma-flat-finite-presentation-CM-fibres-relative-dimension}
または
Morphisms, Lemma \ref{morphisms-lemma-smooth-omega-finite-locally-free}
のいずれからも従うことに注意すれば十分である
% P12-ERR-0166: Japanese closes the frozen parenthetical explanation.
（いずれもスキームについて結果を与え，そこから代数空間についての同様の結果は
Morphisms of Spaces, Lemma
\ref{spaces-morphisms-lemma-openness-bounded-dimension-fibres}
とまったく同様に導かれる．

\medskip\noindent
一般の場合に戻り，始域がスキームである滑らかな全射
$h:V \to \mathcal{T}$ を選ぶ．$v \in V$ ならば，本質的には定義により
$$
\dim_{h(v)}(\mathcal{T}_{f(h(v))}) =
\dim_{v}(V_{f(h(v))}) - \dim_{v}(V_{h(v)}).
$$
$V$ はスキームなので，この等式の右辺の第一項が上半連続であること
（さらに $f$ が滑らかならば局所定数であること）は既に証明した．
一方，第二項は実際に局所定数である．したがって，これらの差は上半連続であり
（$f$ が滑らかならば局所定数であり），よって関数
$\dim_{h(v)}(\mathcal{T}_{f(h(v))})$ は $|V|$ 上で上半連続である
（$f$ が滑らかならば局所定数である）．
射 $|V| \to |\mathcal{T}|$ は開かつ全射なので，補題が従う．
\end{proof}

\noindent
議論を先へ進める前に，スキームの次元論に関する二つの補題を証明する．

\medskip\noindent
最初の補題の背景を説明する．$X$ が有限次元スキームならば，$\dim X$ は
次元 $\dim_x X$ の上限として定義されるので，
点 $x \in X$ であって $\dim_x X = \dim X$ を満たすものが存在する．
次の補題は，さらに点 $x$ を有限型に取れることを示す．

\begin{lemma}
\label{stacks-geometry-lemma-dimension-achieved-by-finite-type-point}
$X$ が有限次元スキームならば，閉点（したがって有限型の点）$x \in X$ であって
$\dim_x X = \dim X$ を満たすものが存在する．
\end{lemma}

\begin{proof}
$d = \dim X$ とおき，$X$ の既約閉部分集合からなる極大な真の降鎖を選ぶ．
これを
\begin{equation}
\label{stacks-geometry-equation-maximal-chain}
Z_0 \supset Z_1 \supset \ldots \supset Z_d.
\end{equation}
と書く．部分集合 $Z_d$ は $X$ の極小な既約閉部分集合であり，
したがって $Z_d$ の任意の点は $Z_d$ の一般点である．
スキーム $X$ の台位相空間は sober なので，$Z_d$ は
ただ一つの閉点 $x \in X$ からなる一点集合であると結論できる．
$U$ を $x$ の任意の近傍とすると，鎖
$$
U\cap Z_0 \supset U\cap Z_1 \supset \ldots \supset U\cap Z_d = Z_d =
\{x\}
$$
は $U$ の既約閉部分集合からなる真の降鎖である．したがって
$\dim U \geq d$ である．よって $\dim_x X \geq d$ を得る．
逆の不等式は明らかなので，補題が証明された．
\end{proof}

\noindent
次の補題は，$\dim_x X$ が，いくつかの穏やかな追加仮定を満たす
既約スキーム上で {\it 定数}関数であることを示す．

\begin{lemma}
\label{stacks-geometry-lemma-constancy-of-dimension}
$X$ を既約，Jacobson，鎖状，かつ局所 Noether 的な有限次元スキームとする．
このとき，$X$ の任意の空でない開部分集合 $U$ に対して
$\dim U = \dim X$ が成り立つ．同値な言い方をすれば，
$\dim_x X$ は $X$ 上の定数関数である．
\end{lemma}

\begin{proof}
二つの主張の同値性は定義から直ちに従う．そこで，$U\subset X$ を
空でない開部分集合とする．明らかに $\dim U \leq \dim X$ であるから，
$\dim U \geq \dim X.$ を示さなければならない．
$d = \dim X$ とおき，$X$ の既約閉部分集合からなる極大な真の降鎖を選ぶ．
これを
$$
X = Z_0 \supset Z_1 \supset \ldots \supset Z_d.
$$
と書く．$X$ は Jacobson なので，極小な既約閉部分集合 $Z_d$ は
ある閉点 $x$ に対する $\{x\}$ に等しい．

\medskip\noindent
$x \in U,$ ならば
$$
U = U \cap Z_0  \supset U\cap Z_1 \supset \ldots \supset
U\cap Z_d = \{x\}
$$
は $U$ の既約閉部分集合からなる真の降鎖である．したがって
$\dim U \geq d$ となり，所望の結果を得る．よって，
$x \not\in U.$ と仮定してよい．

\medskip\noindent
平坦射 $\Spec \mathcal{O}_{X,x} \to X$ を考える．
空でない（したがって稠密な）開部分集合 $U$ の $X$ 上での引き戻しは，
開部分集合 $V \subset \Spec \mathcal{O}_{X,x}$ である．
$U$ を空でない準コンパクトな，したがって Noether 的な開部分集合に
取り替えることにより，包含射 $U \to X$ は準コンパクト射であると仮定してよい．
% P12-ERR-0168: Japanese supplies the citation transition omitted in the frozen prose.
準コンパクト射のスキーム論的像の形成は平坦基底変換と可換であることが
Morphisms, Lemma
\ref{morphisms-lemma-flat-base-change-scheme-theoretic-image}
から分かるので，$V$ は $\Spec \mathcal{O}_{X,x}$ で稠密であり，
特に空でない．もちろん $x \not\in V.$ である
（ここでは $x$ を $\Spec \mathcal{O}_{X,x}$ の閉点を表すためにも用いている．
その像は与えられた点 $x \in X$ に等しいからである）．
さて，$\Spec \mathcal{O}_{X,x} \setminus \{x\}$ は
Properties, Lemma
\ref{properties-lemma-complement-closed-point-Jacobson}
により Jacobson である．したがって，$V$ は，ある閉点 $z$ を含み，これは
$\Spec \mathcal{O}_{X,x} \setminus \{x\}$ の閉点である．
$X$ における $z$ の像の閉包は，ある既約閉部分集合 $Z$ である．これは
$X$ の部分集合で $x$ を含み，$U$ との共通部分は空でない．さらに，$Z$ に真に含まれ，かつ
$\{x\}$ を真に含む既約閉部分集合は存在しない
（$\Spec \mathcal{O}_{X,x}$ への引き戻しは，$X$ の $x$ を含む既約閉部分集合と
$\Spec \mathcal{O}_{X,x}$ の既約閉部分集合との間の全単射を誘導するからである）．
$U \cap Z$ は $U$ の空でない閉部分集合なので，点 $u$ であって $X$ で閉であるものを含む
（$X$ は Jacobson だからである）．また，$U\cap Z$ は既約集合 $Z$ の
空でない（したがって稠密な）開部分集合であり
（この既約集合は $U$ に属さない点 $x$ を含む），包含 $\{u\} \subset U\cap Z$ は真である．

\medskip\noindent
$X$ は鎖状なので，鎖
$$
X = Z_0 \supset Z \supset \{x\} = Z_d
$$
は長さ $d+1$ の鎖に細分でき，それは必ず
$$
X = Z_0 \supset W_1 \supset \ldots \supset W_{d-1} = Z \supset \{x\} = Z_d.
$$
という形でなければならない．$U\cap Z$ は空でないので，
$$
U = U \cap Z_0 \supset U \cap W_1\supset \ldots \supset U\cap W_{d-1}
= U\cap Z \supset \{u\}
$$
は $U$ の既約閉部分集合からなる長さ $d+1$ の真の降鎖である．したがって
$\dim U \geq d$ となり，所望の結果を得る．
\end{proof}
% END EXPANDED MODULE ch107_03.tex
% BEGIN EXPANDED MODULE ch107_04.tex

\noindent
以下の Lemma \ref{stacks-geometry-lemma-constancy-of-dimension} のスタック論的な類似を，
Lemma \ref{stacks-geometry-lemma-irreducible-implies-equidimensional} において証明する．
しかしその前に，追加の定義を導入しなければならない．これは，
% P12-ERR-0169
スキームが鎖状であるという概念が \'etale 局所的なものではないという事実に
よるものである
（Algebra, Remark
\ref{algebra-remark-universally-catenary-does-not-descend} の例を参照されたい）．
この事実のため，代数空間または代数スタックが鎖状であることの意味を
定義するのは困難である
（\cite[page 3]{Osserman} の議論を参照されたい）．
次の定義は，次元論のある側面に対して，欠けている
鎖状代数スタックという概念のよい代用物を与えるように思われる．

\begin{definition}
\label{stacks-geometry-definition-pseudo-catenary}
局所 Noether な代数スタック $\mathcal{X}$ が
{\it 擬鎖状}であるとは，始域が普遍鎖状スキームである
滑らかな全射 $U \to \mathcal{X}$ が存在することをいう．
\end{definition}

\begin{example}
\label{stacks-geometry-example-pseudo-catenary}
$\mathcal{X}$ が普遍鎖状な局所 Noether スキーム $S$ 上
局所有限型であり，$U\to \mathcal{X}$ がスキームを始域とする
滑らかな全射であるとする．このとき合成
% P12-ERR-0170
$U \to \mathcal{X} \to S$ は局所有限型であり，したがって
Morphisms, Lemma
\ref{morphisms-lemma-universally-catenary-local}
により $U$ は普遍鎖状である．
ゆえに $\mathcal{X}$ は擬鎖状である．
\end{example}

\noindent
次の補題は，擬鎖状であるという性質が有限型射を通じて
移ることを示す．

\begin{lemma}
\label{stacks-geometry-lemma-catenary-covers}
% P12-ERR-0171
$\mathcal{X}$ を擬鎖状な局所 Noether 代数スタックとし，
$\mathcal{Y} \to \mathcal{X}$ を局所有限型射とする．
このとき，始域が普遍鎖状スキームである滑らかな全射
$V \to \mathcal{Y}$ が存在する．したがって
$\mathcal{Y}$ も擬鎖状である．
\end{lemma}

\begin{proof}
仮定により，始域が普遍鎖状スキームである滑らかな全射
$U \to \mathcal{X}$ を見つけることができる．
このとき基底変換 $U\times_{\mathcal{X}} \mathcal{Y}$ は
代数スタックである．$V \to U\times_{\mathcal{X}} \mathcal{Y}$ を，
始域がスキームである滑らかな全射とする．
合成 $V \to U\times_{\mathcal{X}} \mathcal{Y} \to \mathcal{Y}$ は
滑らかな全射である（滑らかな全射の合成だからである）．一方，射
$V \to U\times_{\mathcal{X}}
\mathcal{Y} \to U$ は局所有限型である
（局所有限型射の合成だからである）．$U$ は普遍鎖状なので，
Morphisms, Lemma
\ref{morphisms-lemma-universally-catenary-local}
により $V$ は普遍鎖状である．以上で主張が従う．
\end{proof}

\noindent
ここで，関数 $\dim_x(\mathcal{X})$ の $|\mathcal{X}|$ 上での振る舞いを，
ある局所 Noether スタック $\mathcal{X}$ について，
$|\mathcal{X}|$ の既約成分に関して調べる．
併せて，いくつかの関連する話題も扱う．

\begin{lemma}
\label{stacks-geometry-lemma-irreducible-implies-equidimensional}
$\mathcal{X}$ を Jacobson かつ擬鎖状な局所 Noether 代数スタックとし，
$|\mathcal{X}|$ は既約であるとする．
このとき $\dim_x(\mathcal{X})$ は $|\mathcal{X}|$ 上の定数関数である．
\end{lemma}

\begin{proof}
$\dim_x(\mathcal{X})$ が $|\mathcal{X}|$ 上局所定数であることを
示せば十分である．実際，そのときこれは定数となる
（$|\mathcal{X}|$ は既約，したがって連結だからである）．
$\mathcal{X}$ は擬鎖状なので，滑らかな全射 $U \to \mathcal{X}$ であって，
$U$ が普遍鎖状スキームとなるものを見つけることができる．
$\{U_i\}$ を $U$ の準コンパクト開部分スキームによる被覆とすると，
$U$ を $\coprod U_i,$，で置き換えることができ，
関数 $u \mapsto \dim_{f(u)}(\mathcal{X})$ が $U_i$ 上局所定数であることを
示せば十分である．これを一度に一つの $U_i$ について確認するので，
以下では添字を省き，単に $U$ と書いて $U_i$ とは書かない．
$U$ は準コンパクトなので，有限個の既約成分の和，
例えば $T_1 \cup \ldots \cup T_n$ である．各 $T_i$ は Jacobson，
鎖状，かつ局所 Noether であることに注意する．実際これは，
Jacobson，鎖状，かつ局所 Noether なスキーム $U$ の閉部分スキームだからである．

\medskip\noindent
Lemma \ref{stacks-geometry-lemma-behaviour-of-dimensions-wrt-smooth-morphisms} により，
$\dim_{f(u)}(\mathcal{X}) = \dim_{u}(U) - \dim_{u}(U_{f(u)}).$
Lemma \ref{stacks-geometry-lemma-relative-dimension-is-semi-continuous} (2) により，
右辺の第二項は $U$ 上局所定数である．これは $f$ が滑らかだからである．
したがって，$\dim_u(U)$ が $U$ 上局所定数であることを示さなければならない．
$\dim_u(U)$ は，次元 $\dim_u T_i$ の最大値である．ここで $T_i$ は
$U$ の成分のうち $u$ を含むものを動く．したがって，点 $u$ が二つの異なる成分，
例えば $T_i$ と $T_j$（ただし $i \neq j$）に属するとき，
$\dim_u T_i = \dim_u T_j$ であることを示せば十分である．
そのうえで，関数 $t\mapsto \dim_t T$ が，既約な Jacobson，鎖状，かつ
局所 Noether なスキーム $T$ 上で定数であることに注意すればよい
（Lemma \ref{stacks-geometry-lemma-constancy-of-dimension} から従う）．

\medskip\noindent
$V = T_i \setminus (\bigcup_{i' \neq i} T_{i'})$
および $W = T_j \setminus (\bigcup_{i' \neq j} T_{i'})$ とおく．
$V$ と $W$ はいずれも $U$ の空でない開部分集合なので，
$|\mathcal{X}|$ におけるそれぞれの像も空でない開部分集合である．
$|\mathcal{X}|$ は既約なので，$|\mathcal{X}|$ のこれら二つの
空でない開部分集合は空でない交わりをもつ．
$x$ をこの交わりに属する点とし，$v \in V$ と
$w\in W$ を $x$ に写る点とする．すると
$$
\dim T_i = \dim V = \dim_v (U) = \dim_x (\mathcal{X}) + \dim_v (U_x)
$$
であり，同様に
$$
\dim T_j = \dim W = \dim_w (U) = \dim_x (\mathcal{X}) + \dim_w (U_x).
$$
を得る．$u \mapsto \dim_u (U_{f(u)})$ は $U$ 上局所定数であり，
$T_i \cup T_j$ は連結である
（空でない交わりをもつ二つの既約，したがって連結な集合の和だからである）．
ゆえに $\dim_v (U_x) = \dim_w(U_x)$ であり，
直前の二つの等式を比較すると $\dim T_i = \dim T_j$ を得る．
これが求める結論である．
\end{proof}

\begin{lemma}
\label{stacks-geometry-lemma-closed-immersions}
$\mathcal{Z} \hookrightarrow \mathcal{X}$ を
局所 Noether 代数スタックの閉埋め込みとする．
$z \in |\mathcal{Z}|$ の像を $x \in |\mathcal{X}|$ とすると，
$\dim_z (\mathcal{Z}) \leq \dim_x(\mathcal{X})$ である．
\end{lemma}

\begin{proof}
始域がスキームである滑らかな全射
$U\to \mathcal{X}$ を選ぶ．基底変換された射
$V = U\times_{\mathcal{X}} \mathcal{Z} \to \mathcal{Z}$ も
滑らかな全射であり，射影 $V \to U$ は閉埋め込みである．
$v \in |V|$ が $z \in |\mathcal{Z}|$ に写り，
$u$ を $v$ の $|U|$ における像とすると，明らかに
$\dim_v(V) \leq \dim_u(U)$ である．一方，
$\dim_v (V_z) = \dim_u(U_x)$ である．これは
Lemma \ref{stacks-geometry-lemma-base-change-invariance-of-relative-dimension} による．したがって
$$
\dim_z(\mathcal{Z})  = \dim_v(V) - \dim_v(V_z)
\leq \dim_u(U) - \dim_u(U_x) = \dim_x(\mathcal{X}),
$$
となり，主張が従う．
\end{proof}

\begin{lemma}
\label{stacks-geometry-lemma-dimension-via-components}
$\mathcal{X}$ を局所 Noether 代数スタックとし，
$x \in |\mathcal{X}|$ とする．このとき
$\dim_x(\mathcal{X}) = \sup_{\mathcal{T}} \{ \dim_x(\mathcal{T}) \} $
である．ここで $\mathcal{T}$ は，$|\mathcal{X}|$ の既約成分のうち
$x$ を通るものすべて（誘導された被約構造を入れたもの）を動く．
\end{lemma}

\begin{proof}
Lemma \ref{stacks-geometry-lemma-closed-immersions} により，
$\dim_x (\mathcal{T}) \leq \dim_x(\mathcal{X})$ である．ここで
$\mathcal{T}$ は点 $x$ を通る各既約成分を動く．
したがって補題を証明するには，
\begin{equation}
\label{stacks-geometry-equation-desired-inequality}
\dim_x(\mathcal{X}) \leq
\sup_{\mathcal{T}} \{\dim_x(\mathcal{T})\}.
\end{equation}
$U\to\mathcal{X}$ をスキームによる滑らかな被覆とする．
$T$ が $U$ の既約成分であるとき，$\mathcal{T}$ をその像の
$\mathcal{X}$ における閉包とする．これは $\mathcal{X}$ の既約成分である．
$u \in U$ を $x$ に写る点とする．このとき
$\dim_x(\mathcal{X})=\dim_uU-\dim_uU_x=\sup_T\dim_uT-\dim_uU_x$
である．ここで上限は，$U$ の既約成分のうち $u$ を通るものを動く．
上限が達成される成分 $T$ を選ぶと，
$\dim_x(\mathcal{T})=\dim_uT-\dim_u T_x$ であることに注意する．
所望の不等式 (\ref{stacks-geometry-equation-desired-inequality}) は，明らかな不等式
$\dim_u T_x \leq \dim_u U_x.$ から従う．
（$\Spec k \to \mathcal{X}$ が $x$ の代表であるとき，
$T\times_{\mathcal{X}} \Spec k$ は
$U\times_{\mathcal{X}}
\Spec k$ の閉部分空間であることに注意する．）
\end{proof}

\begin{lemma}
\label{stacks-geometry-lemma-dimension-at-finite-type-point}
$\mathcal{X}$ を局所 Noether 代数スタックとし，
$x \in |\mathcal{X}|$ とする．このとき，
任意の開部分スタック $\mathcal{V}$ であって，$\mathcal{X}$ の部分であり
$x$ を含むものに対して，
有限型点 $x_0 \in |\mathcal{V}|$ であって
$\dim_{x_0}(\mathcal{X}) = \dim_x(\mathcal{V})$ を満たすものが存在する．
\end{lemma}

\begin{proof}
始域がスキームである滑らかな全射
$f:U \to \mathcal{X}$ を選び，関数
$u \mapsto \dim_{f(u)}(\mathcal{X});$ を考える．
射 $|U| \to |\mathcal{X}|$ は $f$ により誘導され，開写像
（$f$ が滑らかだからである）かつ全射（仮定による）であり，
有限型点を有限型点へ写す
（これは $|\mathcal{X}|$ の有限型点の定義そのものである）．
したがって，任意の $u \in U$ と $u$ の任意の開近傍に対して，
この近傍に属する有限型点 $u_0$ であって
$\dim_{f(u_0)}(\mathcal{X}) = \dim_{f(u)}(\mathcal{X}).$
を満たすものが存在することを示せば十分である．
このように問題を言い換えると，$f$ の全射性はもはや必要でない．
そこで $U$ を問題の点 $u$ の開近傍で置き換えることにより，
各 $u \in U$ に対して有限型点 $u_0 \in U$ であって
$\dim_{f(u_0)}(\mathcal{X}) = \dim_{f(u)}(\mathcal{X}).$
を満たすものが存在することを示す問題に帰着できる．
Lemma \ref{stacks-geometry-lemma-behaviour-of-dimensions-wrt-smooth-morphisms} により
$\dim_{f(u)}(\mathcal{X}) = \dim_u(U) - \dim_u(U_{f(u)}),$
である一方，
$\dim_{f(u_0)}(\mathcal{X}) = \dim_{u_0}(U) - \dim_{u_0}(U_{f(u_0)}).$
である．$f$ は滑らかなので，式 $\dim_{u_0}(U_{f(u_0)})$ は，
$u_0$ が $U$ 上を動くとき局所定数である
（Lemma \ref{stacks-geometry-lemma-relative-dimension-is-semi-continuous} (2) による）．
したがって必要ならば $U$ を $u$ の周りでさらに縮小し，
この式は定数であると仮定できる．ゆえに問題は，有限型点
$u_0 \in U$ であって $\dim_{u_0}(U) = \dim_u(U)$ を満たすものを
見つけられることを示すことになる．
定義により $\dim_u U$ は，次元 $\dim V$ の最小値である．ここで $V$ は，
開近傍 $V$ のうち $u$ を含む $U$ のものを動く．そこで $U$ を $u$ の周りでさらに縮小し，
$\dim_u U = \dim U$ と仮定できる．
所望の点 $u_0$ の存在は
Lemma \ref{stacks-geometry-lemma-dimension-achieved-by-finite-type-point} から従う．
\end{proof}

\begin{lemma}
\label{stacks-geometry-lemma-monomorphing-a-component-in-of-the-right-dimension}
% P12-ERR-0172
$\mathcal{T} \hookrightarrow \mathcal{X}$ を代数スタックの
局所有限型なモノ射とし，$\mathcal{X}$（したがって $\mathcal{T}$ も）
Jacobson，擬鎖状，かつ局所 Noether であるとする．
さらに $\mathcal{T}$ はある（有限な）次元 $d$ の既約スタックであり，
$\mathcal{X}$ は被約で，次元が $d$ 以下であると仮定する．
このとき空でない開部分スタック $\mathcal{V}$ であって，$\mathcal{T}$ の部分であり，
誘導されるモノ射 $\mathcal{V} \hookrightarrow \mathcal{X}$ が開埋め込みであり，
$\mathcal{V}$ を $\mathcal{X}$ のある既約成分の開部分集合と同一視するものが
存在する．
\end{lemma}

\begin{proof}
$f:U \to \mathcal{X}$ を始域がスキームである滑らかな全射とする．
$\mathcal{X}$ が被約なので，このスキームも必然的に被約である．
$U' = \mathcal{T}\times_{\mathcal{X}} U$ と書く．基底変換された射
% P12-ERR-0173
$U' \to U$ は代数空間の局所有限型なモノ射であり，したがって
Morphisms of Spaces, Lemma
\ref{spaces-morphisms-lemma-locally-quasi-finite-separated-representable} および
\ref{spaces-morphisms-lemma-monomorphism-loc-finite-type-loc-quasi-finite}
により表現可能である．$U$ はスキームなので $U'$ もスキームである．
射影 $f': U' \to \mathcal{T}$ も滑らかな全射である．
$u' \in U'$ とし，その像を $u \in U$ とする．
Lemma \ref{stacks-geometry-lemma-base-change-invariance-of-relative-dimension} により
$\dim_{u'}(U'_{f(u')}) = \dim_u(U_{f(u)}),$
である．一方，$\dim_{f'(u')}(\mathcal{T}) =d
\geq \dim_{f(u)}(\mathcal{X})$ である．これは
Lemma \ref{stacks-geometry-lemma-irreducible-implies-equidimensional} および
$\mathcal{T}$ と $\mathcal{X}$ に関する仮定による．したがって
\begin{equation}
\label{stacks-geometry-equation-dim-inequality}
\dim_{u'} (U') = \dim_{u'} (U'_{f(u')}) + \dim_{f'(u')}(\mathcal{T})
\\
\geq \dim_u (U_{f(u)}) + \dim_{f(u)}(\mathcal{X}) = \dim_u (U).
\end{equation}
$U' \to U$ は局所有限型なモノ射なので，特に非分岐である．
したがって，非分岐射の \'etale 局所構造
\'Etale Morphisms, Lemma \ref{etale-lemma-finite-unramified-etale-local}
により，可換図式
$$
\xymatrix{
V' \ar[r]\ar[d] & V \ar[d] \\
U' \ar[r] & U
}
$$
であって，スキーム $V'$ は空でなく，縦の矢印は \'etale であり，
上の水平矢印は閉埋め込みであるものを見つけることができる．
$V$ を，その像が $U'$ の像と空でない交わりをもつ
準コンパクト開部分集合で置き換え，$V'$ を $V$ の逆像で置き換えることにより，
$V$（したがって $V'$ も）が準コンパクトであるとさらに仮定できる．
$V$ は局所 Noether でもあるので Noether であり，
有限個の既約成分の和である．

\medskip\noindent
% P12-ERR-0174
\'Etale 射は各点での次元を保つので
Descent, Lemma \ref{descent-lemma-dimension-at-point-local}，
(\ref{stacks-geometry-equation-dim-inequality}) から次が従う：
任意の点 $v' \in V'$ と，その像 $v \in V$ に対して，
$\dim_{v'}( V') \geq \dim_v(V)$ である．
特に $V'$ の像は，$V$ の二つの異なる既約成分の交わりに
含まれることはできない．したがって $V$ の既約な開部分集合であって
$V'$ と空でない交わりをもつものを少なくとも一つ見つけることができる．
$V$ をこの部分集合で置き換えることにより，$V$ は整であると仮定できる
（被約かつ既約だからである）．直前の次元の不等式から，
閉埋め込み $V' \hookrightarrow V$ は実際には同型であると結論する．
$W$ を $V'$ の $U'$ における像とすると，$W$ は
$U'$ の空でない開部分集合である（\'etale 射は開写像だからである）．
誘導されるモノ射 $W \to U$ は \'etale である
（これは始域上 \'etale 局所的に，すなわち\ $V'$ へ引き戻した後に成り立つ）．
したがってこれは開埋め込みである（\'etale モノ射だからである）．
$\mathcal{V}$ を $W$ の $\mathcal{T}$ における像とすると，
$\mathcal{V}$ は $\mathcal{T}$ の稠密な（同値なことに，空でない）
開部分スタックであり，その像は $\mathcal{X}$ のある既約成分で稠密である．
最後に，
% P12-ERR-0175
射 $\mathcal{V} \to \mathcal{X}$ は滑らかであることに注意する
（滑らかな射 $W\to \mathcal{V}$ との合成が滑らかだからである）．
またこれはモノ射でもあり，したがって開埋め込みである．
\end{proof}

\begin{lemma}
\label{stacks-geometry-lemma-dims-of-images}
% P12-ERR-0176
$f: \mathcal{T} \to \mathcal{X}$ を，Jacobson，擬鎖状，かつ
局所 Noether な代数スタックの局所有限型射とする．
始域は既約で，終域は準分離であるとし，
$\mathcal{Z} \hookrightarrow \mathcal{X}$ を $\mathcal{T}$ の
スキーム論的像とする．このとき，すべての
% P12-ERR-0177
$t \in |T|$ に対して
$\dim_t( \mathcal{T}_{f(t)}) \geq \dim \mathcal{T}  - \dim \mathcal{Z}$
であり，等号が成り立つ $|\mathcal{T}|$ の空でない
（同値なことに，稠密な）開部分集合が存在する．
\end{lemma}

\begin{proof}
$\mathcal{X}$ を $\mathcal{Z}$ で置き換えることにより，
$f$ はスキーム論的に支配的であり，かつ $\mathcal{X}$ は既約であると
仮定してよく，実際そう仮定する．
ファイバー次元の上半連続性
（Lemma \ref{stacks-geometry-lemma-relative-dimension-is-semi-continuous} (1)）により，
等式 $\dim_t( \mathcal{T}_{f(t)}) =\dim \mathcal{T}  - \dim \mathcal{Z}$
が，$t$ が $\mathcal{T}$ のある空でない開部分スタックに属するとき
成り立つことを証明すれば十分である．
このため，議論の中では常に $\mathcal{T}$ を空でない
開部分スタックで置き換えてよい．

\medskip\noindent
$T' \to \mathcal{T}$ を始域がスキームである滑らかな全射とし，
$T$ を $T'$ の空でない準コンパクト開部分集合とする．
% P12-ERR-0178
$\mathcal{Y}$ は準分離なので，$T \to  \mathcal{Y}$ は準コンパクトである
（Morphisms of Stacks, Lemma
\ref{stacks-morphisms-lemma-quasi-compact-permanence} を，射
$T \to \mathcal{Y} \to \Spec \mathbf{Z}$ に適用する）．
したがって，$\mathcal{T}$ を $T$ の $\mathcal{T}$ における像で
置き換えるならば，
% P12-ERR-0179
Morphisms of Stacks, Lemma
\ref{stacks-morphisms-lemma-surjection-from-quasi-compact}
に訴えることにより，射 $f:\mathcal{T} \to \mathcal{X}$ は
準コンパクトであると仮定できる．

\medskip\noindent
$U \to \mathcal{X}$ を $U$ がスキームであるような滑らかな全射とする．
このとき Lemma \ref{stacks-geometry-lemma-map-of-components} により，
$U$ の既約な開部分集合 $V$ であって
$V \to \mathcal{X}$ が滑らかかつスキーム論的に支配的であるものを
見つけることができる．準コンパクト射がスキーム論的に支配的であるという性質は
平坦基底変換で保たれるので，スキーム論的に支配的な射 $f$ の基底変換
$\mathcal{T} \times_{\mathcal{X}} V \to V$ も
スキーム論的に支配的である．このファイバー積への滑らかな全射をもつ
スキームを $Z$ とする．すると
$Z \to \mathcal{T} \times_{\mathcal{X}} V \to V$ も
スキーム論的に支配的である．したがって，$V$ をスキーム論的に支配する
$Z$ の既約成分 $C$ を見つけることができる．
合成 $Z \to \mathcal{T}\times_{\mathcal{X}} V \to \mathcal{T}$ は滑らかであり，
$\mathcal{T}$ は既約なので，Lemma \ref{stacks-geometry-lemma-map-of-components} により
始域の任意の既約成分の像は $|\mathcal{T}|$ で稠密である．
ここで $C$ を，$Z$ の他のすべての既約成分と交わらない
空でない開部分集合 $W$ で置き換える．次に $\mathcal{T}$ と $\mathcal{X}$ を
$W$ と $V$ の像で置き換える
（Lemma \ref{stacks-geometry-lemma-irreducible-implies-equidimensional} を適用すれば，
これによって $\mathcal{T}$ と $\mathcal{X}$ のいずれの次元も
変わらないことが分かる）．
$\mathcal{W}$ を射 $W \to \mathcal{T}\times_{\mathcal{X}} V$ の像とすると，
$\mathcal{W}$ は $\mathcal{T}\times_{\mathcal{X}} V$ の開部分であり
（射 $W \to \mathcal{T}\times_{\mathcal{X}} V$ は滑らかだからである），
既約である（既約スキームの像だからである）．
こうして，可換図式
$$
\xymatrix{
W \ar[dr] \ar[r]  & \mathcal{W} \ar[r] \ar[d] & V \ar[d] \\
& \mathcal{T} \ar[r] & \mathcal{X}
}
$$
を得る．ここで $W$ と $V$ はスキームであり，縦の矢印は滑らかな全射である．
% P12-ERR-0180
斜めの矢印と左上の水平矢印は滑らかであり，誘導される射
$\mathcal{W} \to \mathcal{T}\times_{\mathcal{X}} V$ は開埋め込みである．
この図式を，補題の主張に現れるさまざまな次元の定義と併せて用いることにより，
確認すべきことを，既知であるスキームの場合に帰着する．

\medskip\noindent
$w \in |W|$ を固定し，その像を順に
$w' \in |\mathcal{W}|$，$t \in |\mathcal{T}|$，$v$ in $|V|$，
$x$ in $|\mathcal{X}|$ とする．
本質的には定義から
（$\mathcal{W}$ が $\mathcal{T}\times_{\mathcal{X}} V$ の開部分であり，
基底変換のファイバーがファイバーの基底変換であることを用いる），等式
$$
\dim_v V_x = \dim_{w'} \mathcal{W}_t
$$
および
$$
\dim_t \mathcal{T}_x = \dim_{w'} \mathcal{W}_v.
$$
を得る．Lemma \ref{stacks-geometry-lemma-behaviour-of-dimensions-wrt-smooth-morphisms}
により
（図式の斜めの矢印と右の縦矢印は，$W$ と $V$ をそれぞれ
スタック $\mathcal{T}$ と $\mathcal{X}$ のスキームによる
滑らかな被覆として実現する），
$$
\dim_t \mathcal{T} = \dim_w W - \dim_w W_t
$$
および
$$
\dim_x \mathcal{X} = \dim_v V - \dim_v V_x.
$$
を得る．これらの等式を組み合わせると
$$
\dim_t \mathcal{T}_x - \dim_t \mathcal{T} + \dim_x \mathcal{X}
= \dim_{w'} \mathcal{W}_v - \dim_w W + \dim_w W_t + \dim_v V -
\dim_{w'} \mathcal{W}_t
$$
となる．$W \to \mathcal{W}$ は滑らかな全射なので，
射 $\Spec \kappa(v) \to V$ によって基底変換した後も同じことが成り立つ
（$W \to \mathcal{W}$ を $V$ 上の射と考える）．この滑らかな射から，
次の二つの等式のうち最初のものを得る：
$$
\dim_w W_v - \dim_{w'} \mathcal{W}_v = \dim_w (W_v)_{w'} = \dim_w W_{w'};
$$
第二の等式は，関係する二つのファイバーを直接比較すれば従う．
同様に，$W \to \mathcal{W}$ を $\mathcal{T}$ 上のスキームの射と考え，
点 $t \in |\mathcal{T}|$ のある代表によって基底変換すると，等式
$$
\dim_w W_t - \dim_{w'} \mathcal{W}_t = \dim_w (W_t)_{w'} = \dim_w W_{w'}.
$$
を得る．すべてを合わせると
$$
\dim_t \mathcal{T}_x - \dim_t \mathcal{T} + \dim_x \mathcal{X}
=  \dim_w W_v - \dim_w W + \dim_v V.
$$
を得る．目標は，この等式の左辺が
% P12-ERR-0181
$t$ の空でない開部分集合上で消えることを示すことである．
$w$ が $W$ の空でない開部分集合を動くとき，その像
$t \in |\mathcal{T}|$ は $|\mathcal{T}|$ の空でない開部分集合を動く
（$W \to \mathcal{T}$ は滑らかだからである）．

\medskip\noindent
したがって，次を示すことに帰着された：
$W\to V$ が，既約な局所 Noether スキームの
局所有限型かつスキーム論的に支配的な射であるならば，
点 $w\in W$ の空でない開部分集合であって
$\dim_w W_v =\dim_w W - \dim_v V$ が成り立つものが存在する．
ここで $v$ は $w$ の $V$ における像を表す．
これは標準的な事実であり，読者の便宜のため証明を再掲する．

\medskip\noindent
% P12-ERR-0182
この等式が成り立つかどうかを変えることなく，
$W$ と $V$ をそれぞれの被約部分スキームで置き換えることができる．
したがって，実際これらは整スキームであると仮定できる．
$\dim_w W_v$ は $W,$ 上局所定数なので，必要ならば $W$ を
空でない開部分集合で置き換え，$\dim_w W_v$ は定数，例えば $d$ に
等しいと仮定できる．この開部分集合をアフィンに選ぶことにより，
射 $W\to V$ は実際に有限型であるとも仮定できる．
必要ならば $V$ を空でない開部分集合で置き換え
（その後，この開部分集合上へ $W$ を引き戻す．得られる引き戻しは空でない．
実際，準コンパクトかつスキーム論的に支配的な射の平坦基底変換は，
やはりスキーム論的に支配的だからである），
さらに $W$ は $V$ 上平坦であると仮定できる．
したがって射 $W\to V$ は，
Morphisms, Definition
\ref{morphisms-definition-relative-dimension-d}
の意味で相対次元 $d$ である．Morphisms, Lemma
\ref{morphisms-lemma-rel-dimension-dimension}
から $\dim_w(W) = \dim_v(V) + d,$ が従い，これが求める等式である．
\end{proof}

\begin{remark}
\label{stacks-geometry-remark-negative-dimension}
直前の補題の状況では，必ずしも
$\dim \mathcal{T} \geq \dim \mathcal{Z}$ とは限らないことに注意する．
これは補題の主張にある不等式と矛盾しない．なぜなら，射 $f$ のファイバーも
代数スタックであり，したがって負の次元をもちうるからである．
これは，$k$ を体とし，射
$[\Spec k/\mathbf{G}_m] \to \Spec k$ に補題を適用することで例示される．

\medskip\noindent
補題の主張における射 $f$ が準 DM であると仮定する
（Morphisms of Stacks, Definition
\ref{stacks-morphisms-definition-separated} の意味である．例えば\ 
代数空間によって表現可能な射は準 DM である）．このとき，終域の点上の
射のファイバーは準 DM 代数スタックであり，したがって非負の次元をもつ．
この場合，補題から実際に
$\dim \mathcal{T} \geq \dim \mathcal{Z}$ が従う．
実際，次のより一般的な結果を得る．
\end{remark}

\begin{lemma}
\label{stacks-geometry-lemma-dims-of-images-two}
% P12-ERR-0183
$f: \mathcal{T} \to \mathcal{X}$ を，Jacobson，擬鎖状，かつ
局所 Noether な代数スタックの局所有限型射であって準 DM であるものとする．
始域は既約で，終域は準分離であるとし，
$\mathcal{Z} \hookrightarrow \mathcal{X}$ を $\mathcal{T}$ の
スキーム論的像とする．このとき
$\dim \mathcal{Z} \leq \dim \mathcal{T}$ であり，
さらに次の二条件のうちちょうど一方が成り立つ：
\begin{enumerate}
\item 任意の有限型点
% P12-ERR-0184
$t \in |T|,$ に対して
$\dim_t(\mathcal{T}_{f(t)}) > 0,$ である．この場合
$\dim \mathcal{Z} < \dim \mathcal{T}$ である．または，
\item $\mathcal{T}$ と $\mathcal{Z}$ は同じ次元をもつ．
\end{enumerate}
\end{lemma}

\begin{proof}
直前の Remark で述べたように，準 DM スタックの次元は常に非負である．
したがって，すべての $t \in |\mathcal{T}|$ に対して
$\dim_t \mathcal{T}_{f(t)} \geq 0$ であり，等式
$$
\dim_t \mathcal{T}_{f(t)} = \dim_t \mathcal{T} - \dim_{f(t)} \mathcal{Z}
$$
が，点 $t\in |\mathcal{T}|$ の稠密な開部分集合上で成り立つ．
\end{proof}
% END EXPANDED MODULE ch107_04.tex
% BEGIN EXPANDED MODULE ch107_05.tex





\section{局所環の次元}
\label{stacks-geometry-section-dimension-local-ring}

\noindent
代数スタックは通常の意味では局所環を実際にもたないが，
局所環の次元を次のように定義できる．

\begin{lemma}
\label{stacks-geometry-lemma-dimension-local-ring-pre}
$\mathcal{X}$ を局所 Noether 的代数スタックとする．
$U \to \mathcal{X}$ を平滑射とし，$u \in U$ とする．このとき
$$
\dim(\mathcal{O}_{U, \overline{u}}) -
\dim(\mathcal{O}_{R_u, e(\overline{u})}) =
2\dim(\mathcal{O}_{U, \overline{u}}) -
\dim(\mathcal{O}_{R, e(\overline{u})})
$$
である．ここで $R = U \times_\mathcal{X} U$，射影を
$s, t : R \to U$，対角射を $e : U \to R$ とし，$R_u$ は
$u$ 上の $s : R \to U$ のファイバーである．
\end{lemma}

\begin{proof}
これは，
$s : \mathcal{O}_{U, \overline{u}} \to \mathcal{O}_{R, e(\overline{u})}$
が Noether 的局所環の平坦な局所準同型であり，したがって
$$
\dim(\mathcal{O}_{R, e(\overline{u})}) =
\dim(\mathcal{O}_{U, \overline{u}}) +
\dim(\mathcal{O}_{R_u, e(\overline{u})})
$$
が Algebra, Lemma
\ref{algebra-lemma-dimension-base-fibre-equals-total} により
成り立つことから従う．
\end{proof}

\begin{lemma}
\label{stacks-geometry-lemma-dimension-local-ring}
$\mathcal{X}$ を局所 Noether 的代数スタックとする．
$x \in |\mathcal{X}|$ を有限型点とする
% P12-ERR-0185: Japanese supplies balanced punctuation around the frozen unmatched citation.
（Morphisms of Stacks, Definition
\ref{stacks-morphisms-definition-finite-type-point}）．
$d \in \mathbf{Z}$ とする．次は同値である：
\begin{enumerate}
\item スキーム $U$，平滑射 $U \to \mathcal{X}$，および $x$ に写る
有限型点 $u \in U$ が存在し，
$2\dim(\mathcal{O}_{U, \overline{u}}) -
\dim(\mathcal{O}_{R, e(\overline{u})}) = d$ となる；
\item 任意のスキーム $U$，平滑射 $U \to \mathcal{X}$，および $x$ に写る
有限型点 $u \in U$ に対して，
$2\dim(\mathcal{O}_{U, \overline{u}}) -
\dim(\mathcal{O}_{R, e(\overline{u})}) = d$ となる．
\end{enumerate}
ここで $R = U \times_\mathcal{X} U$，射影を $s, t : R \to U$，
対角射を $e : U \to R$ とし，$R_u$ は $u$ 上の
$s : R \to U$ のファイバーである．
\end{lemma}

\begin{proof}
$x$ の二つの平滑近傍 $(U, u)$ と $(U', u')$ をとり，
$u$ と $u'$ は有限型点であるとする．$U$ と $U'$ を縮小した後，
$u$ と $u'$ は閉点であると仮定してよい
（有限型点の定義による）．そこで全射 \'etale 射
$W \to U \times_\mathcal{X} U'$ を選ぶ．
$W_u$ を $u$ 上の $W \to U$ のファイバーとし，
$W_{u'}$ を $u'$ 上の $W \to U'$ のファイバーとする．
$u$ と $u'$ は $|\mathcal{X}|$ の同じ点に写るから，
$W_u \cap W_{u'}$ は空でない．したがって，$u$ と $u'$ の双方に
写る閉点 $w \in W$ を選べる．これにより，次段落の議論へ帰着する．

\medskip\noindent
$(U', u') \to (U, u)$ を，$u$ と $u'$ が閉点である
$x$ の平滑近傍どうしの平滑射と仮定する．目標は，$(U, u)$ に対して
定義された不変量が $(U', u')$ に対して定義された不変量と同じであると
示すことである．このため，
$\mathcal{O}_{U, u} \to \mathcal{O}_{U', u'}$ は Noether 的局所環の
平坦な局所準同型であり，したがって
$$
\dim(\mathcal{O}_{U', \overline{u}'}) =
\dim(\mathcal{O}_{U, \overline{u}}) +
\dim(\mathcal{O}_{U'_u, \overline{u}'})
$$
が Algebra, Lemma
\ref{algebra-lemma-dimension-base-fibre-equals-total} により成り立つことに
注意する．
（局所環とその狭義 Hensel 化の性質を関係づけるすべての段階は省略する；
More on Algebra, Section
\ref{more-algebra-section-permanence-henselization} を参照せよ．）
一方，
$$
R' = U' \times_{U, t} R \times_{s, U} U'
$$
である．したがって，
$$
\dim(\mathcal{O}_{R', e(\overline{u}')}) =
\dim(\mathcal{O}_{R, e(\overline{u})}) +
\dim(\mathcal{O}_{U'_u \times_u U'_u, (\overline{u}', \overline{u}')})
$$
であることが分かる．補題を証明するには，
$$
\dim(\mathcal{O}_{U'_u \times_u U'_u, (\overline{u}', \overline{u}')}) =
2\dim(\mathcal{O}_{U'_u, \overline{u}'})
$$
を示せば十分である．これは常に真であるとは限らない
（例えば $U'_u$ が曲線であり，$u'$ がこの曲線の生成点である場合）．
しかし，$u'$ は $u$ 上局所有限型な代数空間 $U'_u$ の閉点である．
この場合，まず Varieties, Lemma
\ref{varieties-lemma-dimension-product-locally-algebraic} により
$\dim_{(u', u')}(U'_u \times_u U'_u) = 2\dim_{u'}(U'_u)$ であり，
次に局所代数的スキームの閉点における次元が局所環の次元と一致するため，
等式が成り立つ；Varieties, Lemma
\ref{varieties-lemma-dimension-locally-algebraic} を参照せよ．
これらのスキームに対する結果を代数空間の言葉に移す作業は省略する．
\end{proof}

\begin{definition}
\label{stacks-geometry-definition-dimension-local-ring}
$\mathcal{X}$ を局所 Noether 的代数スタックとする．
$x \in |\mathcal{X}|$ を有限型点とする．
{\it $x$ における $\mathcal{X}$ の局所環の次元}が
$d \in \mathbf{Z}$ であるとは，Lemma
\ref{stacks-geometry-lemma-dimension-local-ring} の同値な条件が満たされることをいう．
\end{definition}

\noindent
これは確かに Lemma \ref{stacks-geometry-lemma-dimension-local-ring-pre} および
Properties of Stacks, Definition
\ref{stacks-properties-definition-dimension-at-point} によって動機づけられる．
% P12-ERR-0186: Japanese preserves the frozen distinction between the announced dim_x quantity and the following local-ring formula.
本節の最後に，$\mathcal{X}$ の $x$ における半普遍変形環の性質によって
$\dim_x(\mathcal{X})$ を計算する公式を確立する．

\begin{lemma}
\label{stacks-geometry-lemma-dimension-formula}
$\mathcal{X}$ を，局所 Noether 的スキーム $S$ 上局所有限型な
代数スタックとする．$x_0 : \Spec(k) \to \mathcal{X}$ を射とし，
$k$ は $S$ 上有限型な体とする．Remark \ref{stacks-geometry-remark-groupoid-defo}
のように，$\mathcal{F}_{\mathcal{X}, k, x_0}$ を，剰余体 $k$ をもつ
Noether 的完備局所 $S$-代数の余亜群 $(A, B, s, t, c)$ によって表示する．
このとき
$$
\text{局所環の次元：}\mathcal{X}\text{ の }x_0 =
2\dim A - \dim B
$$
である．
\end{lemma}

\begin{proof}
$s \in S$ を $x_0$ の像とする．$\mathcal{O}_{S, s}$ が G-環ならば
（これは実際にはほとんど常に満たされる条件である），
補題を次のように証明できる．Lemma
\ref{stacks-geometry-lemma-Artin-approximation-by-smooth-morphism} により，
% P12-ERR-0187: Japanese supplies the article absent from the frozen English clause.
始域がスキームであり，剰余体 $k$ の点 $u_0 \in U$ を含む平滑射
$U \to \mathcal{X}$ であって，誘導される射
$\Spec(k) \to U \to \mathcal{X}$ が $x_0$ と一致し，かつ
$A = \mathcal{O}_{U, u_0}^\wedge$ となるものを見いだせる．
$R = U \times_\mathcal{X} U$ と書く．このとき
$\mathcal{O}_{R, e(u_0)}^\wedge$ を $B$ と同一視できる．
したがって等式は定義から従う．

\medskip\noindent
証明の残りでは一般の場合の証明法を説明するが，
読者にはこの部分を飛ばすことを勧める．

\medskip\noindent
まず，右辺が $(A, B, s, t, c)$ の選択に依存しないことを示す．
実際，$(A', B', s', t', c')$ を第二の選択とする．
$A$ と $A'$ は $\mathcal{X}$ の $x_0$ における半普遍変形環なので，
必要なら $A$ と $A'$ を入れ替えた後，与えられた半普遍形式対象
$\xi$ および $\xi'$ と両立する形式的平滑写像 $A \to A'$ を
$A$ および $A'$ 上で選べる．
$\widehat{\mathcal{C}}_\Lambda$ は余積をもち，それらは $\Lambda$ 上の
完備テンソル積で与えられることを思い出す；Formal Deformation Theory,
Lemma \ref{formal-defos-lemma-CLambdahat-coproducts} を参照せよ．
このとき $B$ は，$\xi$ の $A \widehat{\otimes}_\Lambda A$ への
二つの前進像の間の同型の関手をプロ表現する．$B'$ についても同様である．
したがって
$$
B' =
B \otimes_{(A \widehat{\otimes}_\Lambda A)}
(A' \widehat{\otimes}_\Lambda A')
$$
である．
$$
A \widehat{\otimes}_\Lambda A \longrightarrow
A \widehat{\otimes}_\Lambda A' \longrightarrow
A' \widehat{\otimes}_\Lambda A'
$$
が，相対次元を形式的平滑写像 $A \to A'$ の相対次元の $2$ 倍とする
形式的平滑写像であることは直ちに分かる．
（これは一般原理から従うが，この特別な場合には $A'$ が $A$ 上の
$r$ 変数形式的冪級数環であることからも分かる．）
したがって，望みどおり $B \to B'$ は相対次元
$2(\dim(A') - \dim(A))$ の形式的平滑写像である．

\medskip\noindent
% P12-ERR-0188: Japanese starts the frozen lowercase sentence normally.
次に，$l/k$ を有限拡大とする．誘導される点を
$x_{l, 0} : \Spec(l) \to \mathcal{X}$ とする．公式の右辺は，
$x_0$ に対しても $x_{l, 0}$ に対しても同じであると主張する．
これは Lemma \ref{stacks-geometry-lemma-versal-ring-field-extension} のように
$A \to A'$ を選び，直前の段落とまったく同様に議論すれば示せる．
詳細は省略する．

\medskip\noindent
% P12-ERR-0189: Japanese grammar supplies the article omitted by the frozen English clause.
最後に，Lemma \ref{stacks-geometry-lemma-versal-ring-flat} の証明と同様に議論すれば，
前二段落の両立性を用いて，$A$ がある $\mathcal{X}$ 上平滑なスキーム
$U$ の点 $u_0$ における完備局所環であり，$u_0$ が有限型点である場合
（第一段落で論じた場合）へ帰着できる．詳細は省略する．
\end{proof}








% END EXPANDED MODULE ch107_05.tex
